\documentclass[11pt]{article}
\usepackage[T1]{fontenc}
\usepackage[utf8]{inputenc}
\usepackage{lmodern}
\usepackage{amsmath,amssymb,amsthm,mathtools}
\usepackage[margin=1in]{geometry}
\usepackage{longtable,booktabs,array}
\usepackage{xurl}
\usepackage[hidelinks]{hyperref}
\newtheorem{theorem}{Theorem}[section]
\newtheorem{proposition}[theorem]{Proposition}
\newtheorem{lemma}[theorem]{Lemma}
\newtheorem{corollary}[theorem]{Corollary}
\theoremstyle{remark}
\newtheorem{remark}[theorem]{Remark}
\DeclareMathOperator{\Gal}{Gal}
\allowdisplaybreaks
\setlength{\emergencystretch}{3em}
\title{ZETA\_NEXT\\Effective class numbers from quadratic zero exclusion,\\
bad fibres of $\Lambda$-schemes, and the non-CM gate}
\author{GPT-6 Astra Pro\\Research run of 21 November 2026}
\date{21 November 2026}
\begin{document}
\maketitle

\begin{abstract}
Using the supplied theorem of Alp\"oge and the explicitly checked
constants in Stark's exceptional-zero argument, we prove
$\operatorname{Res}_{s=1}\zeta_K(s)>1/(173\log D_K)$ for every normal
$K\ne\mathbb Q$, and $1/(173n!\log D_K)$ for every field of degree
$n\ge2$.  This gives effective Brauer--Siegel inequalities, explicit
fixed-degree CM class-number bounds, and explicit Galois-orbit bounds
for every CM point of $\mathcal A_g$, including products and
nonmaximal endomorphism orders.  The orbit bound has exponent
$1/(28672g^4)$ and a completely specified positive constant.
Only the averaged Colmez formula is needed.

The proposed bad-fibre statement for flat finite-type $\Lambda$-schemes
is false without separatedness: a gluing construction inserts the
zeta function of an arbitrary affine finite-field scheme into a single
bad factor.  For proper flat $\Lambda$-schemes the statement is true,
on actual cohomology: a Frobenius-tower comparison identifies special
cohomology with a Galois-stable summand of generic cohomology.
We also prove the corrected proper $\Lambda$-RH conclusion.

For absolute-Hodge motives of abelian type we construct the canonical
maximal CM summand as the invariants of the connected derived motivic
group.  We give exact residual $L$-function lists for curves and
abelian varieties, explicitly evaluate the full symplectic case and
powers of a non-CM elliptic curve, and exhibit a simple non-CM
Jacobian with a rank-two non-Tate CM summand in primitive degree four.
The splitting theorems are unconditional.  Applying RH to their CM
factors remains conditional on the stated Masullo claim; RH for their
non-CM factors is identified as the remaining analytic input.
\end{abstract}

\tableofcontents
\clearpage
\section*{Reading conventions, source status, and the answer to the task}
\addcontentsline{toc}{section}{Reading conventions and answer to the task}

This paper continues ZETA\_TWO and its intake of 21 November 2026.
Its IA01 audit and \path{CONDITIONAL_AUDIT.tsv} stand as requested;
we do not repeat them.  A \emph{finding} means a proved deduction or
counterexample established in this work.  It does not assert novelty
against all of the literature.  Every finding below has received a
complete proof reading; the methods, actors, dates, and limited scope
of each check appear in \path{CHECK_EVENTS.tsv}.

\paragraph{The supplied zero theorem A.}
We cite as a theorem the supplied four-page version of L.~Alp\"oge,
\emph{Remark} \cite{Alpoge}.  The endpoint used here is step~20 on
p.~4, not step~17.  In \path{ALPOGE_Main.lean} the declarations are
\path{riemannHypothesis}, \path{riemannZeta_eq_zero_re},
\path{LFunction_ne_zero_of_half_lt_re},
\path{LFunction_eq_zero_re}, and
\path{LFunction_eq_zero_re_of_isPrimitive}; their precise predicates
are in \path{ALPOGE_Statements.lean}.  We use RH for $\zeta$ and
Dirichlet GRH, especially nonvanishing to the right of $1/2$ for
quadratic characters.  We do not infer GRH for a general Dedekind
zeta function from this input.  No step of the theta, positivity,
or Lean argument is re-audited.

\paragraph{The separate conditional input M.}
The Masullo input is the draft's claim in \S1.2, PDF p.~12, for all
unitary Hecke characters of every number field, including
infinite-order characters.  That assertion is unnumbered in the
supplied draft.  Theorem~13.3.1.2, PDF p.~522, is its intended
\emph{Divisor of $L(s,\chi)$} endpoint but has no completed statement
or proof; \S14.1.2, PDF p.~528, discusses the intended consequence.
These are the exact source statuses already recorded in ZETA\_TWO
MDOC01 \cite{Masullo,ZetaTwo}.  Every use of M in this paper is
explicitly conditional.  Anderson's same-weight factorization,
Corollary~5.7.2 and Remark~5.7.3, p.~181 \cite{Anderson1986}, is
carried through ZETA\_TWO CM01; no new proof of that known theorem
is claimed.

\paragraph{Principal answers.}
For a number field write $D_K=|\operatorname{disc}K|$,
$w_K=|\mu(K)|$ and $(r_1,r_2)$ for its signature.  NX01--NX03 prove
\[
 h_KR_K>
 \frac{w_K\sqrt{D_K}}
 {173\,2^{r_1}(2\pi)^{r_2}\log D_K}
 \quad(K/\mathbb Q\ \hbox{normal},\ K\ne\mathbb Q),
\]
and give the explicit factorial loss for arbitrary degree.
NC01--NC04 propagate this to CM class numbers, heights, and all CM
orbits.  In particular, for an explicit $c_g>0$ and all CM fields
$E$ of degree $2g$,
\[
 h_E>c_gD_E^{1/4}(\log D_E)^{-g}.
\]
The full orbit formula, with its constants defined before and in
NC04, is
\[
 |\Gal(\overline{\mathbb Q}/\mathbb Q)\cdot x|
 \ge \frac{|\operatorname{disc}Z(\operatorname{End}_{\overline{\mathbb Q}}A_x)|^{1/P_g}}
 {s_gT_g^{1/P_g}},\qquad P_g=28672g^4.
\]
There is no residual ineffective constant in these displayed
statements.  They do not themselves supply an algorithm for every
other step of a general Andr\'e--Oort proof.

NX-L01 refutes the exact bad-fibre conjecture asked in the prompt:
its hypotheses do not include separatedness.  NX-L02--NX-L03 prove
the proper-flat splitting assertion, and NX-L04 proves that the
multiplicities need not be the generic inertia multiplicities.
NX-L05 explains a failed proper counterexample route; NX-L06 gives
meromorphic continuation and half-integer divisor lines in the
proper case using A.

NG01 closes ZT06's CM-summand gap for every pure factor of abelian
type.  NG02--NG03 compute the residual lists; NG04--NG07 make their
content concrete.  For a general abelian variety the list is an
exact abstract formula; evaluating its primitive CM projectors
requires the actual Mumford--Tate representation.  No general
algorithm recovering that representation from equations is asserted.
The non-CM RH statements which the question asks us to name remain
the analytic inputs explicitly identified by these formulas.

\paragraph{Finite factors and the meaning of $\zeta_X$.}
For a variety over $\mathbb Q$, the Euler product is understood
through a chosen model $\mathcal X$ over an open of
$\operatorname{Spec}\mathbb Z$; its incomplete form is denoted by
$\zeta_X^S$.  For a finite-type scheme over $\mathbb Z$ it is the
closed-point Euler product.  Restoring bad primes always uses the
actual fibres, as in ZETA\_TWO ZT04.  No bad local factor is silently
replaced by the canonical generic inertia factor.

\paragraph{Archive conventions.}
The available cote-and-scan locators are carried from the supplied
predecessor and intakes, not newly viewed manuscripts.  They are
listed precisely in the final framework table and
\path{IMPLICATIONS.md}.  None of the new analytic estimates,
perfectoid comparisons, or Faltings--Deligne theorems is attributed
to an unavailable Grothendieck scan.  Where the supplied material
does not provide a cote and scan, that absence is stated.

\part{Checked findings}


\section{Explicit residues and effective Brauer--Siegel}
\label{sec:normal}

Throughout, $D_K=|\operatorname{disc}(K)|$, $n_K=[K:\mathbf Q]$,
$(r_1,r_2)$ is the signature, $w_K=|\mu(K)|$, and
$\kappa_K=\operatorname*{Res}_{s=1}\zeta_K(s)$.
The regulator uses the usual logarithmic embedding with a factor $2$
at complex places, and is $1$ in unit rank zero. Thus
\begin{equation}\label{eq:classformula}
 h_KR_K=\frac{w_K\sqrt{D_K}}{2^{r_1}(2\pi)^{r_2}}\kappa_K.
\end{equation}
This is the analytic class number formula, in precisely the normalization
of Garcia--Lee \cite[(1), p.~2]{GarciaLee}.

\subsection{The constants in Stark's reduction}

Here are the inputs, with their original locators checked against the
1974 facsimile. Stark \cite[Lemma~3, pp.~139--140]{Stark} gives at most one
zero in
\[
 \operatorname{Re}s\geq1-\frac1{4\log D_K},\qquad
 |\operatorname{Im}s|\leq\frac1{4\log D_K};
\]
such a zero is real and simple. Write it as $\beta_0$ if it exists, and
set $\beta_1=\beta_0$ in that case, and
$\beta_1=1-(4\log D_K)^{-1}$ otherwise. His Lemma~4, equation~(11),
pp.~140--141, gives
\begin{equation}\label{eq:stark-residue}
 \kappa_K>C_4^{-1}(1-\beta_1),\qquad
 C_4=2\exp\left(\frac{21}{8}
                 +\frac{\gamma+2\log2}{4\log3}\right)<43.2.
\end{equation}
The displayed constant follows by taking $c_1=0$, $c_2=2/\log3$ in
Stark's equation~(7), p.~139, and $c_4=2c_3$ on p.~141.
Garcia--Lee \cite[\S3, pp.~16--17]{GarciaLee} independently records this
same extraction. In particular $4C_4<173$.
For an exact numerical check of the last strict inequality, one may use
$\gamma<289/500$, $\log2<347/500$, and $\log3>549/500$;
the exponent in $C_4/2$ is then less than $3073/1000$.
The rational Taylor-series certificate supplied with this paper proves
$8\exp(3073/1000)<173$ without floating-point comparisons.

Stark's Lemma~8, p.~145, has exactly the constants $4$ for a normal
extension of $\mathbf Q$ and $4n_K!$ otherwise: a real zero in the
corresponding interval descends to a quadratic subfield. Theorem~$1'$,
p.~148, equations~(24)--(25), is the resulting residue statement;
its no-quadratic-subfield clause uses the very same zero exclusion.
We extend that clause using Alp\"oge's theorem, not the other branch of
Theorem~$1'$ involving the constant $c_8$. No numerical assertion about
Stark's tentative value $c_8=\pi/6$ is needed here.

\begin{theorem}[NX01: normal fields, with an absolute constant]
\label{NX01}
Assume the supplied theorem of Alp\"oge in the sense fixed in the
introduction. For every normal number field $K/\mathbf Q$ of degree
$n_K\geq2$,
\begin{equation}\label{eq:normal-lower}
 \boxed{\quad
 \kappa_K>\frac1{173\log D_K},\qquad
 h_KR_K>
 \frac{w_K\sqrt{D_K}}
 {173\,2^{r_1}(2\pi)^{r_2}\log D_K}.
 \quad}
\end{equation}
More precisely, $173$ may be replaced by $4C_4$ with $C_4$ as in
\eqref{eq:stark-residue}. The field $\mathbf Q$ is handled separately:
$D=1$, $\kappa=hR=1$.
\end{theorem}

\begin{proof}
Suppose that a real zero $\beta$ lies in
$[1-(4\log D_K)^{-1},1)$. Since $D_K\geq3$, this interval lies strictly
to the right of $1/2$. Stark's Lemma~8 gives a quadratic subfield
$F\subset K$ with $\zeta_F(\beta)=0$. The factorization
$\zeta_F(s)=\zeta(s)L(s,\chi_{D_F})$ contradicts the supplied
Riemann and Dirichlet nonvanishing theorems at $\beta$.
Thus $\beta_0$ does not exist, and \eqref{eq:stark-residue} gives
$\kappa_K>(4C_4\log D_K)^{-1}>(173\log D_K)^{-1}$.
Equation~\eqref{eq:classformula} gives the second inequality.
The conclusion uses only zero exclusion near $1$ for normal Dedekind
zeta functions; it asserts no full GRH for them.
\end{proof}

\begin{theorem}[NX02: transfer through the Galois closure]
\label{NX02}
For every number field $K$ of degree $n\geq2$, under the same input,
\begin{equation}\label{eq:all-fields-lower}
 \boxed{\quad
 \kappa_K>\frac1{173n!\log D_K},\qquad
 h_KR_K>
 \frac{w_K\sqrt{D_K}}
 {173n!\,2^{r_1}(2\pi)^{r_2}\log D_K}.
 \quad}
\end{equation}
The factor $n!$ can be replaced by $1$ for normal $K$. In particular,
for a CM field $E$ of degree $2g$,
\begin{equation}\label{eq:cm-hR}
 h_ER_E>\frac{w_E\sqrt{D_E}}
 {173(2g)!(2\pi)^g\log D_E}.
\end{equation}
\end{theorem}

\begin{proof}
For clarity we spell out the normal-closure mechanism behind Stark's
Lemma~8, rather than treating a lower bound for the residue of the
closure as a lower bound for the residue of its subfield.
Let $L$ be the Galois closure, $m=[L:\mathbf Q]\leq n!$, and express $L$
as the compositum of $a\leq n$ conjugates of $K$. Stark's Lemma~7,
p.~143, gives
\[
 D_L\mid D_K^{am/n},\qquad \log D_L\leq m\log D_K\leq n!\log D_K.
\]
The Aramata--Brauer divisibility used in the proof of Stark's
Lemma~8 says that $\zeta_L/\zeta_K$ is entire, since $L/K$ is normal.
Consequently a real zero of $\zeta_K$ in
$[1-(4n!\log D_K)^{-1},1)$ would also be a zero of $\zeta_L$ in its
Stark interval, contrary to NX01's zero-exclusion argument.

If the exceptional zero $\beta_0$ in Stark's Lemma~4 for $K$ exists,
it therefore satisfies $1-\beta_0>(4n!\log D_K)^{-1}$; if it does not,
$1-\beta_1=(4\log D_K)^{-1}$ is at least this large.
Apply \eqref{eq:stark-residue} and then
\eqref{eq:classformula}. For a CM field $r_1=0$, $r_2=g$.
\end{proof}

\begin{theorem}[NX03: a quantitative Brauer--Siegel statement]
\label{NX03}
Let $K/\mathbf Q$ be normal, $n\geq2$, and put $L=\log D_K$. Then
\begin{align}
 \frac{w_K\sqrt{D_K}}{173\,2^{r_1}(2\pi)^{r_2}L}
 &<h_KR_K\notag\\
 &\leq\frac{w_K\sqrt{D_K}}{2^{r_1}(2\pi)^{r_2}}
       \left(\frac{eL}{2(n-1)}\right)^{n-1}.
 \label{eq:effective-BS}
\end{align}
In particular the following error interval is explicit:
\begin{align*}
 -\frac n2\log(2\pi)-\log(173L)+\log2
 &<\log(h_KR_K)-\frac12L\\
 &\leq\log(2n^2)+(n-1)\log\left(\frac{eL}{2(n-1)}\right).
\end{align*}
Along any sequence of normal fields with $D_K\to\infty$ and
$n/\log D_K\to0$, these inequalities give
$\log(h_KR_K)/\log\sqrt{D_K}\to1$ with the displayed effective error.
For arbitrary fields of fixed degree, the same conclusion and error
hold with the additional term $-\log(n!)$ in the lower bound.
\end{theorem}

\begin{proof}
The lower bound is NX01. The upper bound is Louboutin
\cite[Theorem~1(1), p.~1195]{NC-Louboutin2001}, inserted into the class number
formula. For the error interval use $w_K\geq2$ and
$2^{r_1}(2\pi)^{r_2}\leq(2\pi)^{n/2}$ in the lower estimate.
For the upper estimate, $w_K\leq2n^2$: indeed
$\varphi(w_K)\leq n$ and the prime-power product formula gives
$\varphi(t)^2\geq t/2$ for every positive integer $t$.
Divide the resulting logarithmic bounds by $L$. If $x=(n-1)/L\to0$,
then $x\log(e/(2x))\to0$, proving the assertion, and the other error
terms visibly tend to zero. The fixed-degree assertion uses NX02.
\end{proof}

\subsection{An independent upper bound relevant to the earlier programme}

\begin{proposition}[NX04: a class-number upper bound needs no new RH]
\label{NX04}
For any number field $K$ of degree $n\geq2$,
\begin{equation}\label{eq:minkowski-class-upper}
 h_K\leq\sqrt{D_K}\left(1+\tfrac12\log D_K\right)^{n-1}.
\end{equation}
Thus a previously established sufficient power
$2h_K|D_p|/|\mu(K)|$ can be bounded above effectively by substituting
\eqref{eq:minkowski-class-upper}, whenever the other factors of that
earlier statement have their specified finite meaning.
\end{proposition}

\begin{proof}
Minkowski's ideal-class bound \cite[Theorem~4.3, p.~70]{MilneANT} gives an integral ideal in every class
of norm at most
$M=(4/\pi)^{r_2}(n!/n^n)\sqrt{D_K}\leq\sqrt{D_K}$.
For the last inequality one may use $n!/n^n\leq2^{1-n}$ and
$r_2\leq n/2$; for $n\geq2$ the resulting factor
$2\pi^{-n/2}$ is at most $1$.
The Euler factors of $\zeta_K$ show that its nonnegative Dirichlet
coefficients are bounded coefficientwise by those of $\zeta(s)^n$.
If $d_n(m)$ denotes the latter coefficients, induction on $n$ gives,
for $X\geq1$,
\[
 \sum_{m\leq X}d_n(m)\leq X(1+\log X)^{n-1}.
\]
Indeed the induction step follows by summing
$\sum_{a\leq X}(X/a)(1+\log X)^{n-2}$ and using
$\sum_{a\leq X}1/a\leq1+\log X$.
Count ideals up to $\sqrt{D_K}$ to conclude.
\end{proof}

\paragraph{Geometric and programme interpretation.}
NX01--NX03 concern $\zeta_{\operatorname{Spec}\mathcal O_K}$ and the
arithmetic special-value combination $h_KR_K$; they give quantitative
information about the integral Picard group and the archimedean unit
lattice. They do not produce a decomposition of the corresponding
\'{e}tale complex, nor an integral $K$-theory calculation from zero
location. This distinction is exactly that of ZETA\_TWO ZT08, SV01--SV02.
The theorem extends ZETA\_TWO IA05's zero-exclusion reduction to explicit
residue bounds. NX04 applies to the sufficient-power expression in
SA06 B.4e as recorded in \texttt{INTAKE\_RESEARCH\_SA06\_20261120.md},
\S4.1; it does not repair the different refuted assertion B.4c.
A lower bound for $hR$ would require an \emph{upper} bound for $R$ to
give a lower bound for $h$. The CM argument below avoids that problem
by an exact regulator quotient. An upper bound for $hR$, in contrast,
needs a regulator lower bound to give an upper bound for $h$.

No manuscript statement of these new explicit estimates has been
identified in Grothendieck's archive. The exact programme anchors
available in the supplied record are SA06 B.4a--e and the determinant
discussion of cote~60, scans~3, 5, 7, 9, as recorded in ZETA\_TWO AR01
and its archive section. These anchors are cited as inherited readings,
not as new facsimile inspections or claims that the archive contains
the present analytic theorem.


% Consequences of the effective residue estimate, /root/cm_effectivity.
% IDs NC01--NC04. Uses theorem/proposition/proof environments and amsmath/amssymb.
% Bibliography keys are supplied in cm/bibliography.tex.

\section{Effective CM class numbers and CM-point orbits}

Throughout this section, $D_E$ denotes the positive absolute discriminant of
$E$, and $\kappa_E=\operatorname*{Res}_{s=1}\zeta_E(s)$.
We use the quadratic Dirichlet GRH theorem supplied as Alp\"oge's theorem
and carried in ZETA\_TWO IA05.  No statement in this section uses Masullo's
claim, the general individual Colmez conjecture, or GRH for a nonquadratic
Hecke character.  The residue estimate used from NX02 is
\begin{equation}\label{eq:NC-residue-input}
 \kappa_E>\frac{1}{173[ E:\mathbf Q]!\log D_E}
 \qquad([E:\mathbf Q]\geq2).
\end{equation}
For normal $E$, NX01 permits the factorial to be replaced by $1$.

\begin{theorem}[NC01: explicit Siegel bounds for CM class numbers]
\label{NC01}
Let $E$ be any CM field of degree $2g$, let $F=E^+$, and put
\[
 c_g=\frac{2}{173(2\pi)^g(2g)!}
       \begin{cases}
       1,&g=1,\\
       \bigl(4(g-1)/e\bigr)^{g-1},&g\geq2.
       \end{cases}
\]
Then
\begin{equation}\label{eq:NC01-log}
 h_E\ \geq\ \frac{h_E}{h_F}\ >\
 c_g\frac{D_E^{1/4}}{(\log D_E)^g}.
\end{equation}
Consequently, for every explicitly given $\varepsilon>0$,
\begin{equation}\label{eq:NC01-epsilon}
 h_E>c_g\left(\frac{e\varepsilon}{g}\right)^g
                 D_E^{1/4-\varepsilon}.
\end{equation}
If $E$ is normal, the same statements hold with $c_g$ multiplied by
$(2g)!$.  In particular, for $H\geq1$, the condition $h_E\leq H$ implies
\begin{equation}\label{eq:NC01-cutoff}
 D_E<\left\{\frac{H}{c_g(e/(8g))^g}\right\}^{8}.
\end{equation}
For imaginary quadratic fields one retains the natural, stronger exponent:
\begin{equation}\label{eq:NC01-quadratic}
 h_E>\frac{\sqrt{D_E}}{173\pi\log D_E}
 \geq\frac{e\varepsilon}{173\pi}D_E^{1/2-\varepsilon}
 \qquad(\varepsilon>0).
\end{equation}
The constants are deliberately coarse and are completely explicit.
The stronger large-discriminant quadratic estimates in
ZETA\_TWO IA02--IA03 remain available; they are not reproved here.
\end{theorem}

\begin{proof}
The relative analytic class number formula, in the normalization of
\cite[\S3, equation (21), p.~1200]{NC-Louboutin2001}, is
\begin{equation}\label{eq:NC01-relative}
 \frac{h_E}{h_F}=
 \frac{Q_Ew_E}{(2\pi)^g}
 \sqrt{\frac{D_E}{D_F}}\frac{\kappa_E}{\kappa_F},
 \qquad Q_E\in\{1,2\},\quad w_E\geq2.
\end{equation}
This formula accounts for the regulator ratio
$R_E/R_F=2^{g-1}/Q_E$; replacing a bound for $h_ER_E$ by a bound for
$h_E$ without this step would be invalid.
For $g\geq2$, Louboutin's explicit upper bound
\cite[Theorem~1(1), p.~1195]{NC-Louboutin2001} gives
\[
 \kappa_F\leq
 \left(\frac{e\log D_F}{2(g-1)}\right)^{g-1}.
\]
Combining this with \eqref{eq:NC-residue-input} gives the stronger
field-dependent bound
\begin{equation}\label{eq:NC01-strong}
 \frac{h_E}{h_F}>
 \frac{2}{173(2\pi)^g(2g)!\log D_E}
 \sqrt{\frac{D_E}{D_F}}
 \left(\frac{2(g-1)}{e\log D_F}\right)^{g-1}.
\end{equation}
The discriminant identity
$D_E=D_F^2N_{F/\mathbf Q}(\mathfrak d_{E/F})$ implies
$D_F\leq D_E^{1/2}$.  Hence the square root in
\eqref{eq:NC01-strong} is at least $D_E^{1/4}$, while
$\log D_F\leq\tfrac12\log D_E$.  This proves
\eqref{eq:NC01-log} for $g\geq2$.  For $g=1$ use
$F=\mathbf Q$, $\kappa_F=1$, and
$D_E^{1/2}\geq D_E^{1/4}$ directly in
\eqref{eq:NC01-relative}.  Finally $h_F\geq1$.

For $t>0$, the maximum of $t^g e^{-\varepsilon t}$ occurs at
$t=g/\varepsilon$ and equals $(g/(e\varepsilon))^g$.
With $t=\log D_E$, this proves \eqref{eq:NC01-epsilon}.
Taking $\varepsilon=1/8$ proves \eqref{eq:NC01-cutoff}.
Normality changes only the residue constant, as stated.
For \eqref{eq:NC01-quadratic}, a quadratic field is normal and has
$R_E=1$, so NX01 and the ordinary class number formula give
$h_E=w_E\sqrt{D_E}\kappa_E/(2\pi)>
\sqrt{D_E}/(173\pi\log D_E)$.
The same calculus inequality with $g=1$ gives its second inequality.
\end{proof}

\paragraph{Arithmetic-geometric meaning of NC01.}
Here $h_E=|\operatorname{Pic}(\operatorname{Spec}\mathcal O_E)|$.
The unit logarithm determinant in $K_1(\mathcal O_E)$ is the regulator;
\eqref{eq:NC01-relative} cancels this determinant against its counterpart
for $\operatorname{Spec}\mathcal O_F$.  On the generic fibres,
$\zeta_E/\zeta_F$ is the $L$-function of the genuine Artin summand
$\operatorname{Ind}_{G_F}^{G_{\mathbf Q}}\chi_{E/F}$ in
$H^0(\operatorname{Spec}E)$, where $\chi_{E/F}$ is the quadratic
character of $G_F$.  The passage through the normal closure in NX02
controls a possible exceptional quadratic constituent, not every zero
of this Artin $L$-function.  Thus NC01 changes the effectivity of the
class-number estimate without asserting the full RH of $\zeta_E$.

\begin{theorem}[NC02: an explicit logarithmic CM height bound]
\label{NC02}
For $g\geq1$, define
\[
 a_g=2^{g-1}\left(\frac{75}{2}+16g!\right),
 \qquad b_g=4g(2^g-1),\qquad H_g(D)=a_g\log D+b_g.
\]
Let $A/\overline{\mathbf Q}$ have dimension $g$ and complex
multiplication by the full ring $\mathcal O_E$ of a CM field of degree
$2g$.  Let $h_F$ be the stable height in Faltings' original normalization.
Then
\begin{equation}\label{eq:NC02-height}
 h_F(A)\leq H_g(D_E).
\end{equation}
This assertion holds for every CM type, including imprimitive types.
\end{theorem}

\begin{proof}
We spell out the constants in the proved argument of
\cite[Proposition~5.3 and its proof, pp.~19--20]{NC-Wei2021}.
Use first the height $h_W$ defined with the squared archimedean norm
$(2\pi)^{-g}\int|\omega\wedge\overline\omega|$.
This is the normalization of Yuan--Zhang
\cite[\S1.1, p.~534; Theorem~1.1, p.~535]{NC-YuanZhang2018}
and Wei.  Gaudron--R\'emond
\cite[\S2.3, pp.~351--352]{NC-GRPeriods2014} give
\begin{equation}\label{eq:NC02-normalization}
 h_W(A)=h_F(A)+\frac g2\log\pi.
\end{equation}
Their Theorem~1.1, p.~345, has positive left hand side and upper bound
$16h_F(A)+49g$, so it implies $h_F(A)>-49g/16$.
In particular $h_W(A)>-4g$.

Stark's Lemma~9, pp.~145--146, applied with $[F:\mathbf Q]=g$, says
that a real zero of $L(s,\chi_{E/F})=\zeta_E(s)/\zeta_F(s)$ in
\[
 1-\frac{1}{16g!\log D_E}\leq s<1
\]
would descend to an imaginary quadratic subfield of $E$.
Alp\"oge's quadratic Dirichlet GRH excludes this.  The only possible
zero in the wider square of side $1/(4\log D_E)$ is real and simple
(Stark Lemma~3; also \cite[\S1, p.~3]{NC-Wei2021}), so
$C_{\mathrm{zero}}(g)=16g!$ satisfies the hypothesis of Wei's
Proposition~5.3.

Let $\Phi_0$ be the type of $A$.  The averaged Colmez identity is
\[
 2^{-g}\sum_\Phi h_W(E,\Phi)=
 -\frac12\frac{L'(0,\chi_{E/F})}{L(0,\chi_{E/F})}
 -\frac14\log(D_E/D_F).
\]
Wei's estimate in the proof of Proposition~5.3 bounds the right hand
side above by
\[
 \frac12\left\{
 \left(\frac{75}{2}+16g!\right)\log D_E
 +\frac g2\bigl(\psi(1/2)-\log\pi\bigr)\right\},
\]
where $\psi$ is the logarithmic derivative of the gamma function.
Isolate the summand $\Phi_0$ and use $h_W(E,\Phi)>-4g$ for the
other $2^g-1$ summands.  Since
$\psi(1/2)-\log\pi=-\gamma-\log(4\pi)<0$, discarding the last
term gives $h_W(A)\leq H_g(D_E)$.
Equation \eqref{eq:NC02-normalization} then gives
\eqref{eq:NC02-height}.  No individual Colmez formula was used.
\end{proof}

\paragraph{Geometric meaning of NC02.}
The height is the normalized arithmetic degree of the metrized Hodge
line $\det(e^*\Omega^1_{\mathcal A/\mathcal O_K})$ on an abelian
scheme after a field extension giving good reduction.
Averaging the CM types expresses that degree in terms of the logarithmic
derivative of a weight-zero Artin $L$-function.  The polarizable
weight-one CM realization $H^1(A)$ is present through this Hodge line;
its own infinite-order Hecke $L$-function is not assumed to satisfy GRH.
NC02 strengthens the averaged-height consequence recorded after
ZETA\_TWO IA05 and does not settle the individual Colmez identity
which IA05 explicitly leaves as a separate input.

\subsection{Explicit constants for the orbit statements}

For each integer $r\geq1$, let $c_r,a_r,b_r$ have the values just
specified and define
\begin{align*}
 \alpha_r&=2^{10}r^3,& B_r&=(14r)^{64r^2},&
 \eta_r&=\frac{1}{16\alpha_r},\\
 U_r&=a_r+\frac{b_r+1}{\log3},&
 q_r&=2+\frac{r}{2\alpha_r},&
 s_r&=2\cdot3^{4r^2},\\
 f_r&=\frac{c_r^{1/(2\alpha_r)}}
              {2^{r/\alpha_r}B_rU_r^2},&&&\\
 \ell_r&=\min\left\{1,\,
           f_r\left(\frac{e\eta_r}{q_r}\right)^{q_r}\right\}.&&&
\end{align*}
Every constant here is positive and given by elementary operations.
The isogeny estimate to be used is
\cite[Theorem~1.4, p.~2060; definition of $\kappa(A)$, p.~2058]
{NC-GRPol2014}: two isogenous abelian varieties $A,A'$ of dimension $r$
defined over a number field of degree $k$ have an isogeny in each
direction of degree at most
\begin{equation}\label{eq:NC-GR}
 \bigl(B_rk\max\{h_F(A),\log k,1\}^{2}\bigr)^{\alpha_r}.
\end{equation}
The theorem permits the isogeny over any extension on which the two
varieties are isogenous.  In the applications below there is already
an isogeny over the stated field.

\begin{theorem}[NC03: primitive full-ring CM orbits]
\label{NC03}
Let $A/\overline{\mathbf Q}$ be geometrically simple of dimension $g$
with $\operatorname{End}(A)=\mathcal O_E$ for a CM field $E$ of degree
$2g$.  If $A$ and all its endomorphisms are defined over $K$, then
\begin{equation}\label{eq:NC03-def}
 [K:\mathbf Q]\geq\ell_g D_E^{\eta_g}.
\end{equation}
For a polarized $A$ with field of moduli $\mathbf Q(A)$,
\begin{equation}\label{eq:NC03-moduli}
 [\mathbf Q(A):\mathbf Q]\geq
 \frac{\ell_g}{s_g}D_E^{1/(16384g^3)}.
\end{equation}
This is an effective version of Tsimerman's Theorems~1.2 and~4.2.
\end{theorem}

\begin{proof}
Put $D=D_E$, $k=[K:\mathbf Q]$, and fix the primitive CM type $\Phi$
of $A$.  The usual CM ideal construction gives $h_E$ isomorphism
classes in $S(E,\Phi)$; the quotient by the kernel of an integral ideal
$I\subset\mathcal O_E$ has isogeny degree $N(I)$, and every class
arises in this way.  The precise facts used are recorded in
\cite[\S2.2, p.~382, proof of Proposition~2.2]{NC-Tsimerman2018}.
Because all endomorphisms of $A$ are defined over $K$, their common
ideal kernels and these quotients are defined over $K$ as well.
For each quotient $A_I$, apply \eqref{eq:NC-GR} to $A,A_I$ over $K$.
It gives an isogeny $A\to A_I$ of degree at most
\[
 N=\bigl(B_gk\max\{H_g(D),\log k,1\}^{2}\bigr)^{\alpha_g}.
\]
All homomorphisms in this isogeny class commute with the fixed
$E$-action: after fixing one $E$-linear isogeny, the rational Hom space
is a rank-one module over the commutative algebra
$\operatorname{End}^0(A)=E$.  Thus the isogeny just obtained is among
those classified by integral ideals.  There are consequently at least
$h_E$ integral ideals of norm at most $N$.

For any number field of degree $2g$ and any $N\geq1$,
\[
 \#\{I\subset\mathcal O_E:N(I)\leq N\}
 \leq N^2\zeta_E(2)
 \leq N^2\zeta(2)^{2g}<2^{2g}N^2.
\]
The middle inequality follows by bounding each rational-prime Euler
factor by $(1-p^{-2})^{-2g}$.  Hence $h_E<2^{2g}N^2$.
If $k>D$, \eqref{eq:NC03-def} follows from $\ell_g\leq1$ and
$\eta_g<1$.  Assume therefore $k\leq D$.  As $D\geq3$,
\[
 \max\{H_g(D),\log k,1\}\leq U_g\log D.
\]
The lower bound NC01 and the last ideal count give
\[
 c_g\frac{D^{1/4}}{(\log D)^g}
 <2^{2g}(B_gkU_g^2(\log D)^2)^{2\alpha_g}.
\]
Taking $2\alpha_g$-th roots and solving for $k$ yields
\[
 k>f_g\frac{D^{1/(8\alpha_g)}}{(\log D)^{q_g}}.
\]
For $t>0$, the maximum of $t^{q_g}e^{-\eta_gt}$ is
$(q_g/(e\eta_g))^{q_g}$.  Using $t=\log D$ proves
\eqref{eq:NC03-def} by the definition of $\ell_g$.
Finally Tsimerman's Lemma~4.1, p.~384, gives a field of definition of
$A$ and all its endomorphisms of degree at most
$s_g[\mathbf Q(A):\mathbf Q]$.  (The lemma is formulated with
unpolarized field of moduli and its extension defines every
polarization.  Taking the compositum with the polarized field of
moduli preserves the bound $s_g$ on the relative degree.)
Apply \eqref{eq:NC03-def} to that field.
\end{proof}

\paragraph{Which quantitative steps changed.}
The analytic inputs to the orbit argument are now the explicit
class-number lower bound NC01 and the logarithmic height bound NC02.
The ideal count is elementary.  The isogeny estimate
\eqref{eq:NC-GR} has specified constants, so it leaves no ineffective
constant in NC03.  Its exponent is intentionally sacrificed to keep
one uniform, directly evaluable expression.  This is an orbit estimate;
no numerical degree bound for arbitrary special subvarieties is being
claimed from it.

\begin{theorem}[NC04: all CM points, including nonmaximal orders]
\label{NC04}
For $g\geq1$ define the following further explicit constants:
\begin{align*}
 \beta_g&=16\alpha_g,\\
 M_g&=\max_{1\leq r\leq g}\ell_r^{-16\alpha_r},\\
 A_g^*&=\max_{1\leq r\leq g}a_r,
 &C_g^*&=\max_{1\leq r\leq g}b_r,\\
 V_g&=1+g\{A_g^*(\log M_g+\beta_g)+C_g^*\},\\
 T_g&=M_g^g(B_gV_g^2)^{4g\alpha_g},
 &P_g&=28g\alpha_g=28672g^4.
\end{align*}
Let $x\in\mathcal A_g(\overline{\mathbf Q})$ be a CM point, let
$A_x$ be its principally polarized abelian variety, and let
$R_x=Z(\operatorname{End}_{\overline{\mathbf Q}}A_x)$, an order in a
product of CM fields.  Then
\begin{equation}\label{eq:NC04-orbit}
 |\operatorname{Gal}(\overline{\mathbf Q}/\mathbf Q)\cdot x|
 \geq s_g^{-1}T_g^{-1/P_g}\,|\operatorname{disc}R_x|^{1/P_g}.
\end{equation}
In particular the full Galois-orbit assertion of
\cite[Theorem~5.2, p.~385]{NC-Tsimerman2018} has an explicit exponent
and an explicit multiplicative constant under Alp\"oge's theorem.
\end{theorem}

\begin{proof}
Put $d=[\mathbf Q(x):\mathbf Q]$ and $A=A_x$.
By the same Lemma~4.1 and the compositum just described,
choose $K$ over $\mathbf Q(x)$ with
$k=[K:\mathbf Q]\leq s_gd$ such that $A$ and every geometric
endomorphism are defined over $K$.

R\'emond's maximal-order replacement theorem
\cite[Theorem~1.1, pp.~1--2 of the author preprint]{NC-Remond2017}
provides a $K$-isogenous abelian variety $B/K$ whose endomorphism
ring is a maximal order.  This theorem works over the same base
field: no extension depending on the index of the original order
is needed.  All geometric rational endomorphisms of $B$ are defined
over $K$, by transport through the $K$-isogeny.  The integral ones
are therefore defined over $K$ too.

Repeated application of R\'emond's Theorem~1.2, p.~2, decomposes
$B$ over $K$ as a product of simple factors $B_j$ with maximal
endomorphism rings.  Every geometric endomorphism of each factor is defined over $K$,
by its inclusion into and projection from $B$.  Such a factor is
geometrically simple: applying the same Theorem~1.2 over
$\overline{\mathbf Q}$ to a geometric proper abelian subvariety
would give a nontrivial idempotent endomorphism.  This idempotent
is defined over $K$, contradicting simplicity over $K$.
As $A$ is CM, each factor has dimension $r_j\leq g$,
$\operatorname{End}(B_j)=\mathcal O_{E_j}$ for a CM field of degree
$2r_j$, and primitive CM type (\cite[\S7, paragraph preceding Lemma~7.5,
pp.~17--18 of the preprint]{NC-Remond2017}).  There are at most $g$ factors,
counting repetitions, with $\sum_jr_j=g$.

NC03 in its field-of-definition form implies, for every $j$,
\[
 D_{E_j}\leq(k/\ell_{r_j})^{16\alpha_{r_j}}
          \leq M_g k^{\beta_g}.
\]
The stable Faltings height is additive on products
(\cite[\S2.3, p.~352]{NC-GRPeriods2014}).  NC02 thus gives
\begin{align*}
 h_F(B)&\leq\sum_j\{a_{r_j}\log D_{E_j}+b_{r_j}\}\\
 &\leq g\{A_g^*\log M_g+C_g^*\}
          +gA_g^*\beta_g\log k.
\end{align*}
It follows that
\[
 \max\{h_F(B),\log k,1\}
 \leq V_g(1+\log k)\leq V_gk,
\]
where the last inequality holds for $k\geq1$.
The isogeny estimate \eqref{eq:NC-GR}, applied with $B$ as the
height input, supplies an isogeny $\varphi:B\to A$ of degree
\begin{equation}\label{eq:NC04-degree}
 N\leq(B_gV_g^2)^{\alpha_g}k^{3\alpha_g}.
\end{equation}

Identify the centres of the rational endomorphism algebras by
$\varphi$.  Let $\mathcal O$ be the integral closure of $\mathbf Z$
in this centre.  Since $\operatorname{End}(B)$ is maximal,
its centre is $\mathcal O$: adjoining central algebraic integers
would otherwise give a strictly larger order.  Also
$R_x\subset\mathcal O$, because an integral endomorphism acts on
an integral homology lattice and is integral over $\mathbf Z$.
The kernel of $\varphi$ is killed by its order $N$, so there is an
isogeny $\psi:A\to B$ with $\psi\varphi=[N]_B$.
For $u\in\mathcal O$ the endomorphism
$\varphi u\psi$ of $A$ is central and represents $Nu$ under our
identification.  Thus
\[
 N\mathcal O\subset R_x\subset\mathcal O,
 \qquad [\mathcal O:R_x]\leq N^{\operatorname{rank}_{\mathbf Z}\mathcal O}
                         \leq N^{2g}.
\]
The discriminant of the product order $\mathcal O$ is the product
of the discriminants of its field factors, one for each simple
isogeny component, retaining separate components even if their
fields are isomorphic.  There are at most $g$ of these.  The discriminant-index formula
and \eqref{eq:NC04-degree} give
\begin{align*}
 |\operatorname{disc}R_x|
 &=|\operatorname{disc}\mathcal O|[\mathcal O:R_x]^2\\
 &\leq M_g^gk^{g\beta_g}N^{4g}\\
 &\leq M_g^g(B_gV_g^2)^{4g\alpha_g}
             k^{g\beta_g+12g\alpha_g}
   =T_gk^{P_g}\leq T_g(s_gd)^{P_g}.
\end{align*}
Solving for $d$, which is precisely the Galois-orbit cardinality,
proves \eqref{eq:NC04-orbit}.
\end{proof}

\paragraph{Geometric meaning and exact scope of NC03--NC04.}
The points lie in the moduli scheme $\mathcal A_g/\mathbf Q$.
Isogeny preserves the rational CM motive and all its $\ell$-adic
realizations, while replacing the endomorphism order changes an integral
lattice.  NC04 quantifies that lattice change by
$N\mathcal O\subset R_x$ and the discriminant-index identity.
The Hodge line supplies the height used to bound the isogeny.
Hence these results make the CM-orbit input to Tsimerman's proof
effective, including the centre of a nonmaximal endomorphism ring.
They use the abelian category and its endomorphisms, not a presumed
motivic splitting for arbitrary varieties.  The weight-one non-CM
factors in ZETA\_TWO ZT06 remain a different question.

\paragraph{The distinction from a full effective Andr\'e--Oort algorithm.}
NC04 supplies all constants in the arithmetic orbit lower bound requested
here.  The complete proof in \cite[\S6.3, Theorem~6.1, and \S6.5,
p.~388]{NC-Tsimerman2018} also uses a subpolynomial counting theorem and
an o-minimal reduction from special subvarieties to special points.
Those cited steps, as stated there, supply no algorithm for their
constants or finite families; such an algorithm does not follow from
NC04 alone.  This identifies the scope of the present proof, not a
claim that effective replacements cannot exist.  Indeed, for powers
of the modular curve there already is such a replacement:
Binyamini \cite[\S1.2, Theorem~1, preprint p.~2; \S1.4, pp.~3--4]
{NC-Binyamini2020} gives effective bounds for maximal special
subvarieties once the possible exceptional imaginary quadratic field
is removed by Dirichlet GRH.  That specialization is effective under
the supplied Alp\"oge theorem.  The two nonarithmetic effectivity
issues for the original o-minimal proof are listed separately in the
same \S1.4.

\paragraph{Archive and programme locators.}
These quantitative class-number and moduli-height estimates are not
statements found in EGA, SGA, FGA, or the supplied Grothendieck scans.
Their direct programme targets are ZETA\_TWO IA05 and its intake
\S6(1)(a), and the appended CM class-number discussion in
INTAKE\_RESEARCH\_SA06\_20261120, \S4, rows B.4a--B.4e, restated
in ZETA\_TWO IAQ-008--012.  The supplied intake does not identify a
cote-and-scan for B.4a--B.4e; none is invented here.  The explicit
archive connection specified in ZETA\_TWO's intake \S5.4 is
cote~15, scan~70 for SA06 A.8a, concerning quadratic
class-number/supersingular questions.  The related row A.8c has
neither a separate cote-and-scan locator nor a definition of its
ring $R$ in the supplied intake.  The sharper quadratic estimates
remain ZETA\_TWO IA02--IA03.
The cohomological determinant interpretation is the one already fixed
in ZETA\_TWO ZT08: SGA~6, Exp.~V, \S2 and (7.1.1), and the programme's
cote~60, scans~3--9.  These locators concern the framework, not a
claim that the new numerical inequalities occur in the archive.


% Task 2 contribution. Stable IDs NX-L01--NX-L07.
% Whole-proof readings and their exact scope are recorded in CHECK_EVENTS.tsv.

\section{Bad fibres of $\Lambda$-schemes}

\subsection{The separatedness qualification is substantive}

A $\Lambda$-structure on a flat scheme over $\mathbb Z$ means commuting
endomorphisms $\psi_q$ lifting Frobenius at every prime $q$; the definition
itself does not require separatedness (Borger, introduction, p.~2, and
\S1.1, p.~6).  Borger explicitly adds separatedness at the beginning of
\S6 before Theorem~6.1 and Corollary~6.2.  The conjecture in the present
prompt says ``flat and of finite type'' without adding separatedness.
That literal statement has a counterexample.

\begin{theorem}[NX-L01: arbitrary extra bad factors without separatedness]
\label{NX-L01}
Fix a prime $p$, let $Y=\mathbb A^d_{\mathbb Z}$, and let
$Z\subseteq Y_{\mathbb F_p}$ be any closed subscheme.  Glue two copies
$Y_1,Y_2$ along the identity of $U=Y\setminus |Z|$, and call the resulting
scheme $X$.  Then $X$ is flat and of finite type over $\mathbb Z$ and
admits a $\Lambda$-structure.  Its generic fibre is $Y_{\mathbb Q}$ and
\[
 \zeta_X(s)=\zeta(s-d)\,Z(Z,p^{-s}).
\]
Thus its $p$-factor in the correction of ZETA\_TWO~ZT04 is exactly
$Z(Z,p^{-s})$.  In particular the proposed cyclotomic integral-Tate
factor assertion is false as stated.
\end{theorem}

\begin{proof}
Write $\varphi_q:Y\to Y$ for coordinatewise $q$th powers.  The fold map
$r:X\to Y$ is the identity on each chart, and write $i_1:Y\to X$ for the
first chart.  Because Frobenius on $Y_{\mathbb F_p}$ is the identity on
its topological space, $\varphi_p^{-1}(U)=U$.  Consequently the two
maps $\varphi_p:Y_a\to Y_a$ glue to an endomorphism $\psi_p$ of $X$.
For $q\ne p$, set
\[
 \psi_q=i_1\circ\varphi_q\circ r.
\]
The identities $r i_1=1$, $r\psi_p=\varphi_p r$ and
$\psi_p i_1=i_1\varphi_p$, together with
$\varphi_q\varphi_{q'}=\varphi_{q'}\varphi_q$, prove pairwise
commutativity for all primes.  At $q\ne p$ the overlap is the whole
fibre, so $r:X_{\mathbb F_q}\simeq Y_{\mathbb F_q}$ and $\psi_q$
restricts to its absolute Frobenius.  At $p$ this assertion holds on
both affine charts by construction.  These commuting Frobenius lifts
supply the $\Lambda$-structure by the flat criterion in Borger~\S1.1.
Flatness and local finite type are local on the source; the two affine
charts are flat and of finite type, and their finite union is
quasi-compact.  Thus $X$ satisfies the hypotheses in the assertion.

For each $r\geq1$, the two copies have precisely the points of
$Y(\mathbb F_{p^r})\setminus Z(\mathbb F_{p^r})$ identified.  Hence
\[
 \#X(\mathbb F_{p^r})=p^{dr}+\#Z(\mathbb F_{p^r}).
\]
All other fibres are those of $Y$.  Taking the exponential of
$\sum_{r\geq1}\#X(\mathbb F_{p^r})t^r/r$ proves the local identity,
then the global Euler product and the stated ZT04 correction.

For a completely explicit choice, take $p=5$, $d=2$, and
$Z:\ y^2=x^3-x$ in $\mathbb A^2_{\mathbb F_5}$.  Its smooth projective
completion $E$ (discriminant $64\ne0$ in $\mathbb F_5$) has one
point at infinity and eight points over $\mathbb F_5$: the five affine $x$-values give respectively
$1,1,2,2,1$ choices for $y$.  Silverman, \emph{The Arithmetic of Elliptic Curves}, second edition,
Theorem~V.2.4 (p.~143), gives
\[
 Z(E,t)=\frac{1+2t+5t^2}{(1-t)(1-5t)},\qquad
 Z(Z,t)=\frac{1+2t+5t^2}{1-5t}.
\]
The two numerator inverse roots are $-1+2i,-1-2i$, of absolute value
$\sqrt5$, whereas every inverse root of a product of factors
$(1-\xi5^jt)^{\pm1}$ with $j\in\mathbb Z$ has absolute value an
integral power of $5$.  The numerator has no root in common with
$1-5t$, so cancellation cannot repair the discrepancy.  This proves
refutation.  The chosen $X$ is not separated: the two sections of $X$
obtained by lifting the point $(0,0)$ agree over $\mathbb Z[1/5]$ and
disagree at $5$.  In particular this example says nothing against a
version that expressly assumes separatedness.
\end{proof}

\subsection{A local comparison for a proper Frobenius lift}

For the proper question, a comparison must go beyond proper base
change: proper base change alone does not identify the generic and
special fibres of a singular family.  The following construction uses
the lift at the bad prime.  Its finite-coefficient form is recorded
first because filtered colimits and passage to $\mathbb Q_\ell$ must
not be interchanged without an argument.

\begin{theorem}[NX-L02: the proper Frobenius tower]
\label{NX-L02}
Let $\mathcal X$ be a proper flat scheme of finite type over
$\mathbb Z_p$, and let $\varphi:\mathcal X\to\mathcal X$ be a
$\mathbb Z_p$-endomorphism whose special fibre is absolute Frobenius.
For every prime $\ell\ne p$ and every $a\geq1$, there are natural,
$G_{\mathbb Q_p}$-equivariant isomorphisms
\[
 H^i(\mathcal X_{\overline{\mathbb F}_p},\mathbb Z/\ell^a)
 \simeq
 \mathop{\rm colim}_{\varphi^*}
 H^i(\mathcal X_{\overline{\mathbb Q}_p},\mathbb Z/\ell^a).
 \tag{L2.1}
\]
Here $G_{\mathbb Q_p}$ acts on the left through its unramified
quotient.  In particular
$H^i(\mathcal X_{\overline{\mathbb F}_p},\mathbb Q_\ell)$ is
isomorphic to a $G_{\mathbb Q_p}$-stable direct summand of
$H^i(\mathcal X_{\overline{\mathbb Q}_p},\mathbb Q_\ell)$.
\end{theorem}

\begin{proof}
Put $C=\widehat{\overline{\mathbb Q}_p}$ and let $C^\flat$ be its tilt.
We use the following precise analytic results: Huber (1994),
Remark~4.6(iv)(d), identifies the generic fibre of the formal completion
of a proper scheme with the analytification of its algebraic generic
fibre; Scholze, \emph{Perfectoid spaces}, Theorem~5.2 gives tilting of
perfectoid algebras, Theorem~7.12 gives the equivalence of their
\'etale sites, Definition~7.14 and Corollary~7.18 give cohomological
continuity for perfectoid tilde-limits, and Theorem~9.7 (Huber's
Theorem~3.8.1) gives the algebraic--adic comparison with the finite
coefficients used here.  These are cited, not proved again below.
The construction required to apply them in this setting is as follows.

Complete $\mathcal X$ along its special fibre.  The topological
space of the resulting formal scheme $\mathfrak X$ is
$|\mathcal X_{\mathbb F_p}|$, on which $\varphi$ is the identity.
Every formal affine open $\mathfrak U=\operatorname{Spf}A$ is
therefore invariant under inverse image by $\varphi$, and the map on
it is a continuous endomorphism of $A$ lifting Frobenius modulo $p$.
Set
\[
 B=(A/pA)^{\mathrm{perf}},\qquad
 A_\infty=\widehat{\mathop{\rm colim}_{\varphi} A}^{\,(p)}.
 \tag{L2.2}
\]
The ring $A$ is $p$-torsion-free, by flatness.  Filtered colimits of
$p$-torsion-free $\mathbb Z_p$-modules are $p$-torsion-free.  For any
$p$-torsion-free ring $D$, its $p$-adic completion is again
$p$-torsion-free and its reduction modulo $p^n$ is $D/p^nD$: this
follows by applying the inverse limit to the exact multiplication by
$p$ sequences for $D/p^mD$, whose transition maps are surjective.
Consequently $A_\infty$ is $p$-torsion-free and
$A_\infty/pA_\infty=B$.  Since $B$ is perfect, the strict-$p$-ring
characterization of Witt vectors gives the canonical isomorphism
\[
 A_\infty\simeq W(B).
 \tag{L2.3}
\]
For this characterization see Scholze, Proposition~5.13(i) and
Remark~5.14, or Bhatt, \emph{Lecture II: $\delta$-rings},
Proposition~3.2.  All these identifications commute with restriction
to smaller formal opens.  Indeed, a principal formal open defined by
$f$ has reduction $B_f$ after perfection, and its completed ring is
the unique strict $p$-ring lifting $B_f$; thus the transition maps are
the Witt lifts of the maps on the perfected special fibre.

Define
\[
 D_B=\bigl(W(B)\otimes_{\mathbb Z_p}\mathcal O_C\bigr)^{\wedge}_p,
 \qquad R_B=D_B[1/p].
\]
The $\mathcal O_C$-module $D_B$ is torsion-free: its uncompleted
version is flat over $\mathcal O_C$, $p$-adic completion preserves
$p$-torsion-freeness, and every nonzero scalar of $\mathcal O_C$
divides a power of $p$.  Hence $D_B$ is flat over the valuation ring
$\mathcal O_C$ (also see Scholze, Lemma~5.3(i)).  Its reduction modulo
$p$ is $B\otimes_{\mathbb F_p}\mathcal O_C/p$, and Frobenius induces
an isomorphism $D_B/p^{1/p}\longrightarrow D_B/p$: on the first
factor it is the automorphism of $B$, and on the second it is the
perfectoid Frobenius isomorphism for $\mathcal O_C$.  Thus $D_B^a$
satisfies Definition~5.1(ii) of Scholze, and $R_B$ is perfectoid by
Theorem~5.2.  Under tilting it corresponds to
\[
 R_B^\flat=
 \bigl(B\otimes_{\mathbb F_p}\mathcal O_{C^\flat}\bigr)^{\wedge}[1/\varpi^\flat],
 \tag{L2.4}
\]
where the completion is with respect to any pseudouniformizer.  To
check (L2.4), reduce both integral almost rings modulo matching
pseudouniformizers and use
$\mathcal O_C/p\simeq\mathcal O_{C^\flat}/p^\flat$; their reductions
are the same $B$-algebra.  The uniqueness part of Theorem~5.2 gives
the identification, compatibly with maps in $B$ and with valued-field
automorphisms of $C$.  No choice of a coordinate system or of a
semisimple Frobenius action is involved.

For precision about integral structures, throughout this construction
$\operatorname{Spa}(R_B,R_B^+)$ uses for $R_B^+$ the integral closure
of $D_B$ in $R_B$.  On the characteristic-$p$ side use the integral
closure of $(B\otimes\mathcal O_{C^\flat})^{\wedge}$ in $R_B^\flat$.
These are corresponding plus rings under Scholze's Lemma~6.2: open
integrally closed subrings are determined by their images modulo
the scalar ideal $\mathfrak m_C R_B^\circ$ (respectively
$\mathfrak m_{C^\flat}(R_B^\flat)^\circ$); the two images here
are the integral closures of the same image of
$B\otimes_{\mathbb F_p}\overline{\mathbb F}_p$.  Taking an integral
closure imposes no new condition on valuations already bounded by
the original ring.  Thus (L2.4) identifies the required affinoid
pairs, not just their rings after inverting a pseudouniformizer.

The affinoid perfectoid spaces attached to the $R_B$ glue to a
perfectoid space $\mathcal X_\infty$.  For an explicit check on the
open domains, a principal formal open $D(f)$ has analytic domain
$|f|=1$.  In $W(B)$, the image of $f$ differs from its Teichm\"uller
lift $[\bar f]$ by an element of $pW(B)$, so the same domain is
$|[\bar f]|=1$.  Its rational localization is the completed base
change of $W(B_{\bar f})$.  Finite principal refinements of overlaps
therefore give precisely the same open affinoid domains on both
sides of tilting.  More explicitly, take a finite
formal affine cover of $\mathfrak X$.  The associated affinoid domains
cover $(\mathcal X_C)^{\mathrm{ad}}$ by properness, and their inverse
images under every $\varphi$ are themselves.  On each such domain
$R_B$ is the completion of the direct limit of the affinoid coordinate
rings in the $\varphi$-tower.  A compatible family of bounded
continuous valuations on those rings gives a valuation bounded on
the common integral direct-limit ring.  It is therefore continuous
for the $p$-adic topology and extends uniquely to its completion.  This identifies
the underlying space with the inverse limit of the affinoid spaces;
the induced residue-field maps have dense image.  The same argument
on overlaps glues the identifications.  Hence
\[
 \mathcal X_\infty\ \sim\
 \varprojlim_{\varphi}(\mathcal X_C)^{\mathrm{ad}}
 \tag{L2.5}
\]
in the sense of Scholze's Definition~7.14.

By (L2.4) its tilt is the perfection of the generic fibre of the
constant formal model
$\mathcal X_{\mathbb F_p}\times_{\mathbb F_p}\mathcal O_{C^\flat}$.
For this identification, apply Huber's result first to the
proper noetherian model
$\mathcal X_{\mathbb F_p}\times_{\mathbb F_p}\mathbb F_p[[u]]$,
where $u$ maps to a pseudouniformizer of $C^\flat$, and then base change
to $C^\flat$.  Formal generic fibres commute with that base change on
the displayed affine rings.  This avoids applying Huber's noetherian
statement directly over $\mathcal O_{C^\flat}$.  The generic fibre is
therefore
$(\mathcal X_{\mathbb F_p}\times_{\mathbb F_p}C^\flat)^{\mathrm{ad}}$.
Perfection induces an equivalence of \'etale topoi: use universal
homeomorphism invariance, or Scholze's Corollary~7.19 on the relative
Frobenius tower.  Its finite-coefficient cohomology is therefore that
of $\mathcal X_{\overline{\mathbb F}_p}$, by algebraic--adic comparison
and proper base change (SGA~4, XII, Theorem~5.1 and Corollary~5.2).
Combining this with tilting of \'etale sites and Corollary~7.18 proves
(L2.1).  The construction, reduction, and tilting are functorial for
the semilinear action of $G_{\mathbb Q_p}$ on $C$ and $C^\flat$.
On the constant characteristic-$p$ model the resulting action on
cohomology is exactly the action on the residue field, by proper base
change.  This proves the asserted Galois equivariance.

It remains to justify the rational direct-summand assertion; it is
not an interchange of an inverse limit with a filtered colimit.
Write
$M_i=H^i(\mathcal X_{\overline{\mathbb Q}_p},\mathbb Z_\ell)$,
a finitely generated $\mathbb Z_\ell$-module, and $T=\varphi^*$.
For any finitely generated $\mathbb Z_\ell$-module $M$ with
endomorphism $T$, the sequence $T^{n!}$ converges in
$\operatorname{End}_{\mathbb Z_\ell}(M)$ to an idempotent $e$.
Indeed, on each finite module $M/\ell^aM$ the images of powers of $T$
stabilize, $T$ is a permutation of the stable image, and a sufficiently
large factorial kills the transient part and is divisible by the
order of that permutation.  The resulting idempotents are compatible
as $a$ varies.  Thus
\[
 M=M^{\mathrm{u}}\oplus M^{\mathrm{n}},\quad
 M^{\mathrm{u}}=eM,\quad
 \mathop{\rm colim}_T(M/\ell^aM)=M^{\mathrm{u}}/\ell^aM^{\mathrm{u}}.
 \tag{L2.6}
\]
Every operator commuting with $T$, including the Galois action,
commutes with $e$.  On $M^{\mathrm{u}}$ the map $T$ is invertible,
and on the other summand it is topologically nilpotent.

The universal-coefficient exact sequence is
\[
 0\longrightarrow M_i/\ell^aM_i\longrightarrow
 H^i(\mathcal X_{\overline{\mathbb Q}_p},\mathbb Z/\ell^a)
 \longrightarrow M_{i+1}[\ell^a]\longrightarrow0.
\]
Take the exact filtered colimit under $T$ and use (L2.1) and (L2.6).
The resulting rightmost groups are killed by a fixed power of $\ell$,
independent of $a$, since the torsion subgroup of $M_{i+1}$ is finite.
The leftmost groups are $M_i^{\mathrm{u}}/\ell^aM_i^{\mathrm{u}}$,
with surjective transition maps.  Taking the inverse limit and then
inverting $\ell$ therefore yields
\[
 H^i(\mathcal X_{\overline{\mathbb F}_p},\mathbb Q_\ell)
 \simeq M_i^{\mathrm{u}}\otimes\mathbb Q_\ell
 \subseteq M_i\otimes\mathbb Q_\ell.
 \tag{L2.7}
\]
The inclusion is a $G_{\mathbb Q_p}$-stable direct summand by the
idempotent $e$.  This finishes the proof.
\end{proof}

\begin{theorem}[NX-L03: bad-prime splitting for proper flat $\Lambda$-schemes]
\label{NX-L03}
Let $X$ be proper and flat of finite type over $\mathbb Z$, with a
$\Lambda$-structure.  For every prime $p$ and every $\ell\ne p$,
\[
 H^i(X_{\overline{\mathbb F}_p},\mathbb Q_\ell)
 \simeq\bigoplus_j V_{p,i,j}\otimes\mathbb Q_\ell(-j),
 \tag{L3.1}
\]
where each $V_{p,i,j}$ has finite image under
$G_{\mathbb F_p}$ and is semisimple.  In particular geometric
Frobenius is semisimple and its eigenvalues have the form
$\xi p^j$, with $\xi$ a root of unity and $j\in\mathbb Z$.
After a finite extension of $\mathbb F_p$, every cohomology group is
a direct sum of integral Tate twists.  The bounded derived
$\mathbb Q_\ell$-representation $R\Gamma_c(X_{\overline{\mathbb F}_p},\mathbb Q_\ell)$
likewise splits, noncanonically, into those summands in their respective
degrees.  Consequently every factor of $R_{X,S}$ in ZETA\_TWO~ZT04
is cyclotomic integral-Tate for this proper class, including at bad
primes.
\end{theorem}

\begin{proof}
Apply NX-L02 to $X\times\mathbb Z_p$ and $\varphi=\psi_p$.
Borger's Theorem~6.1 and Corollary~6.2 apply to the proper, hence
separated, scheme $X$: its generic-fibre cohomology is potentially
a direct sum of integral Tate twists, and its global Galois action
is abelian.  Over $\overline{\mathbb Q}_\ell$ it is thus a direct
sum of characters $\epsilon\chi_\ell^{-j}$ with $\epsilon$ of
finite order.  One can see semisimplicity without enlarging the
coefficient field by splitting after the finite-index subgroup in
Corollary~6.2 and averaging an equivariant projection over the finite
quotient.  The $G_{\mathbb Q_p}$-stable direct summand supplied by
NX-L02 is semisimple of the same form.  Its inertia action is trivial
because it is the cohomology of a variety over $\mathbb F_p$.
Therefore only unramified constituents occur, and geometric
Frobenius on each is $\epsilon(F_p)p^j$.  This is (L3.1), with the
finite characters grouped over $\mathbb Q_\ell$.

For the assertion about the derived object, use the elementary
hereditary property of finite-dimensional representations of an
invertible operator over $\mathbb Q_\ell$.  They are torsion modules
over the principal ideal domain $\mathbb Q_\ell[T,T^{-1}]$;
$\operatorname{Ext}^r=0$ for $r\geq2$.  In the continuous subcategory
relevant here the eigenvalues are $\ell$-adic units, and the same
finite-length module calculation applies.  The truncation triangles
of a bounded complex consequently split successively.  This gives
a noncanonical direct sum of its cohomology objects shifted to their
degrees.  Properness identifies compactly supported and ordinary
cohomology.  Notice that this argument does not invoke a motivic
splitting or the Tate conjecture.

Finally the trace formula of SGA~5, XV, \S3, no.~2, unnumbered theorem,
equation~(1), expresses the local zeta function as
\[
 Z(X_{\mathbb F_p},t)=
 \prod_i\det(1-tF_p\mid H^i(X_{\overline{\mathbb F}_p},\mathbb Q_\ell))^{(-1)^{i+1}}.
\]
Every determinant splits into the claimed factors by (L3.1).  The
canonical local factors of the abelian Artin--Tate generic
representations have the same form, since inertia invariants of a
finite-image representation again have finite-image Frobenius.
Taking their ratio is precisely ZT04.  This proves the final claim.
\end{proof}

\begin{proposition}[NX-L04: proper splitting does not preserve multiplicities]
\label{NX-L04}
Even for a finite flat proper $\Lambda$-scheme, the bad special-fibre
cohomology need not equal the generic-fibre cohomology or its inertia
invariants.  In particular NX-L03 does not remove the finite
correction in ZT04.
\end{proposition}

\begin{proof}
Take $X=\mu_p=\operatorname{Spec}\mathbb Z[t]/(t^p-1)$, with
$\psi_q(t)=t^q$.  The quotient is free of rank $p$ over $\mathbb Z$;
the endomorphisms commute, preserve the relation, and lift Frobenius.
The generic geometric fibre has $p$ points, whereas the reduced
special fibre at $p$ is a single rational point.  Thus generic $H^0$
has dimension $p$ and special $H^0$ has dimension $1$.
Moreover $X_{\mathbb Q}$ has one rational point and the cyclotomic
factor $\operatorname{Spec}\mathbb Q(\zeta_p)$ (for $p=2$ the latter
is a second rational point).  Its generic canonical local factor at
$p$ is $(1-t)^{-2}$ because each number-field component has one
residue-degree-one prime over $p$.  Its actual special-fibre local
factor is $(1-t)^{-1}$.  The ZT04 correction is therefore $1-t$.
This also shows directly that special $H^0$ is not the full
inertia-invariant part of generic $H^0$.
\end{proof}

\begin{corollary}[NX-L05: the vertical elliptic blowup has no Frobenius lift]
\label{NX-L05}
Let $E\subset\mathbb P^2_{\mathbb F_p}$ be a smooth plane cubic and set
$\mathcal B=\operatorname{Bl}_{E}(\mathbb P^2_{\mathbb Z_p})$.
This is a proper flat scheme with generic fibre $\mathbb P^2_{\mathbb Q_p}$,
but it has no $\mathbb Z_p$-endomorphism lifting absolute Frobenius on
its special fibre.  In particular a vertical elliptic blowup of a
Tate generic fibre cannot provide a proper counterexample to NX-L03.
More generally, any proper flat model with
$H^i(\mathcal X_{\overline{\mathbb Q}_p},\mathbb Q_\ell)=0$ and
$H^i(\mathcal X_{\overline{\mathbb F}_p},\mathbb Q_\ell)\ne0$ has no
such Frobenius lift.
\end{corollary}

\begin{proof}
The general assertion is immediate from the direct-summand conclusion
of NX-L02.  We compute the indicated nonzero group for the blowup,
so that the obstruction is explicit.  Locally the centre is cut out
by the regular sequence $(p,F)$, with $F\bmod p$ an equation for $E$.
The blowup charts have equations $F=pu$ and $p=Fv$.  The total blowup
is integral, since its Rees algebra embeds in the polynomial ring
of the integral ambient scheme, and is therefore flat over the DVR
$\mathbb Z_p$.  It is projective, and after inverting $p$ the ideal
is the unit ideal, giving the stated generic fibre.

Its reduced special fibre, which here is its whole special fibre,
is the union of the strict transform $D\simeq\mathbb P^2_{\mathbb F_p}$
and the exceptional divisor $P=\mathbb P(N_{E/\mathbb P^2_{\mathbb Z_p}})$,
a $\mathbb P^1$-bundle over $E$; their intersection is the copy of $E$
which is the section corresponding to the rank-one normal bundle
inside the special fibre.  These assertions also follow immediately
from the two displayed blowup charts; the strict transform is unchanged
because blowing up the Cartier divisor $E$ in $\mathbb P^2_{\mathbb F_p}$
is the identity.

The constant-sheaf sequence for the union $D\cup_E P$ is exact on
all geometric stalks.  Its long cohomology sequence contains
\[
 H^2(D)\oplus H^2(P)\longrightarrow H^2(E)
 \longrightarrow H^3(\mathcal B_{\overline{\mathbb F}_p})
 \longrightarrow H^3(D)\oplus H^3(P)\longrightarrow H^3(E).
\]
All coefficients here are $\mathbb Q_\ell$.  The first arrow is
surjective: the pullback of $H^2(E)$ to $P$ restricts identically along
the section $E\hookrightarrow P$.  Also $H^3(D)=H^3(E)=0$.
For the projective-line bundle $P\to E$, proper and smooth base
change give relative cohomology $\mathbb Q_\ell$ in degree zero,
$\mathbb Q_\ell(-1)$ in degree two, and zero otherwise.  The relative
hyperplane class supplies the latter identification.  In its Leray
spectral sequence the only term of total degree three is
$H^1(E,\mathbb Q_\ell)(-1)$, and no differential can reach or leave
it because $E$ has cohomological dimension two.  Consequently
\[
 H^3(\mathcal B_{\overline{\mathbb F}_p},\mathbb Q_\ell)
 \simeq H^1(E_{\overline{\mathbb F}_p},\mathbb Q_\ell)(-1),
\]
a two-dimensional group.  On the generic fibre
$H^3(\mathbb P^2_{\overline{\mathbb Q}_p},\mathbb Q_\ell)=0$.
NX-L02 therefore excludes a Frobenius lift.  In weights language this
is a weight-three elliptic contribution in the special fibre.  The
base-change inputs used in this computation have the exact locators
SGA~4, XII, Theorem~5.1 and Corollary~5.2, and XVI, Theorem~2.1.
\end{proof}

\begin{corollary}[NX-L06: the corrected proper $\Lambda$-RH conclusion]
\label{NX-L06}
Using the stipulated Alp\"oge theorem A, let $X$ be a proper flat
$\Lambda$-scheme of finite type over $\mathbb Z$.  Then $\zeta_X(s)$
has meromorphic continuation, and every non-real zero or pole has
real part in $\tfrac12\mathbb Z$.  More precisely, its global
Dirichlet--Tate factors have their non-real divisors on lines
$\Re(s)=j+\tfrac12$ with $j\in\mathbb Z$, while its actual local
factors and the finite correction factors in ZT04 have their divisors
on lines $\Re(s)=j$ with $j\in\mathbb Z$.
\end{corollary}

\begin{proof}
Borger, Theorem~6.1 and Corollary~6.2, give abelian Artin--Tate
cohomology of $X_{\mathbb Q}$, so on a good open its virtual zeta
factor is a finite product of integral shifts of primitive Dirichlet
$L$-functions with integer exponents.  Restoring the omitted primes
is exactly ZETA\_TWO~ZT04.  NX-L03 makes each actual and correction
factor a product of $(1-\xi p^{j-s})^{\pm1}$, whose divisor has
$\Re(s)=j$.  The supplied Alp\"oge \emph{Remark}, step~20, p.~4,
and the endpoint declarations
\texttt{LFunction\_eq\_zero\_re} and
\texttt{LFunction\_ne\_zero\_of\_half\_lt\_re} in
\texttt{ALPOGE\_Main.lean}, supply the Dirichlet zero-location and
nonvanishing statements.  Together with the classical functional
equations and their real trivial divisors, they place the non-real
divisors of the global shifted factors on $j+\tfrac12$, as in
ZETA\_TWO~ZT01.  Cancellation in a finite meromorphic product removes
divisor support and cannot create new support.  This proves the claim.
No Masullo input enters this corollary.
\end{proof}

\paragraph{Exact scope of the resolution.}
NX-L01 refutes the conjecture literally asked, since separatedness
was not one of its hypotheses.  NX-L02--NX-L03 answer the proper flat
question positively on actual cohomology, including semisimplicity;
NX-L04 shows why equality with the whole generic cohomology or with
its inertia invariants would be a different and false assertion.
A separately proposed \emph{separated, nonproper} strengthening is
not settled by these arguments: the counterexample is nonseparated,
and the formal generic-fibre comparison used in NX-L02 relies on
properness.  No $\Lambda$-equivariant compactification is asserted.
This limitation does not leave the literal conjecture undecided.

\begin{proposition}[NX-L07: cyclotomic local factors do not imply a $\Lambda$-structure]
\label{NX-L07}
The integer ring $\mathbb Z[\zeta_5]$ has cyclotomic integral-Tate
local zeta factors but admits no Frobenius lift at $5$, hence no
$\Lambda$-structure.  Thus the prompt's local-factor observation for
abelian integer rings is valid without asserting that all those
rings are $\Lambda$-rings.
\end{proposition}

\begin{proof}
Every nonzero ideal of the ring of integers $\mathcal O_K$ of a
number field meets $\mathbb Z$ nontrivially: if $0\ne x$ lies in the
ideal, its monic minimal polynomial gives a nonzero integer constant
term in the ideal $(x)$.  Therefore every unital $\mathbb Z$-algebra
endomorphism of $\mathcal O_K$ is injective.  It extends to an
endomorphism of $K$, hence to an automorphism since $K$ is
finite-dimensional over $\mathbb Q$, and preserves its integral
closure $\mathcal O_K$ in both directions.  Its reduction modulo
any integer is consequently an automorphism.

Apply this to $K=\mathbb Q(\zeta_5)$.  Reducing the cyclotomic
polynomial modulo $5$ gives
$\Phi_5(T)=(T-1)^4$ in $\mathbb F_5[T]$, so
\[
 \mathbb Z[\zeta_5]/5\mathbb Z[\zeta_5]
 \simeq\mathbb F_5[\varepsilon]/(\varepsilon^4).
\]
Its Frobenius kills the nonzero element $\varepsilon$, and is not an
automorphism.  It cannot be the reduction of a ring endomorphism.
On the other hand, each local zeta factor of $\operatorname{Spec}\mathcal O_K$
is a product of $(1-t^f)^{-1}$ over the residue degrees $f$ of the
primes above the rational prime.  Each such factor splits into
$(1-\xi t)^{-1}$ over the $f$th roots of unity.  These have integral
Tate index $j=0$, proving the local-factor claim independently of
the nonexistent $\Lambda$-structure.
\end{proof}

\subsection{Placement in Grothendieck's frameworks}

NX-L01 is a statement about flat nonseparated schemes over
$\operatorname{Spec}\mathbb Z$, obtained by Zariski gluing.  Its
local zeta factor is a multiplicative invariant of point counts;
the weight-one factor is the realization of $h^1(E)$ and cannot be
replaced by integral Tate weights.  The full half-integer-line
statement is not refuted: those two inverse roots have weight one,
so their local divisors lie on $\Re(s)=1/2$.

NX-L02 and NX-L03 concern actual $\ell$-adic realizations, including
semisimplicity, and hence strengthen a mere assertion in a
Grothendieck group.  Their local determinants are computed in the
$\lambda$-ring $K_0$ of finite-dimensional representations, using
exterior powers.  This is the formalism of SGA~6, V, \S2 and
\S7, (7.1.1), used already in ZETA\_TWO~ZT08 and SV01.  The sheaf
operations and proper base change live in SGA~4, XII, Theorem~5.1
and Corollary~5.2; compact supports are treated in XVII, \S\S5.1--5.3;
the zeta determinant is SGA~5, XV, \S3, no.~2, equation~(1).
The perfectoid step is a later comparison theorem, not a theorem
silently ascribed to EGA, SGA, or FGA.

For the archive connection, the supplied predecessor places exterior
and symmetric powers and characteristic determinants at cote~60,
scans~3 and~5 (ZETA\_TWO, \emph{IMPLICATIONS}, \S5, placement table);
its scan~7, archival p.~6, records a reciprocal characteristic-polynomial
identity, while scan~9, archival p.~8, records Weil and dual/Tate
normalizations.  These exact locators are carried from the checked
predecessor, not newly inspected scans in this run.  DP\_MULT~DP11
and its intake identify the binomial/identity-Adams special case;
that case does not encompass NX-L01, whose Frobenius lifts are
nonidentity.  DP\_MULT's reported cote~153, PDF p.~3/archive p.~2,
concerns binomial polynomials, not a bad-fibre cohomology theorem.
No motive-level or cycle-level decomposition is asserted by (L3.1).


% Contribution /root/noncm_gate; model GPT-6 Astra Pro; effort not exposed.
% Whole-proof check events are supplied separately in CHECK_EVENTS.tsv.
% Bibliography entries are supplied in bibliography.tex.

\section{The non-CM gate for curves and abelian varieties}
\label{sec:noncm-gate-next}

We use the convention of ZETA\_TWO CM01:
$\operatorname{wt} M(n)=\operatorname{wt}M-2n$ and
$L(M(n),s)=L(M,s+n)$.  ``CM'' in this section means
\emph{potentially CM in the category of motives for absolute Hodge
cycles}, including every Artin motive.  In particular a nonabelian
finite-image Artin representation is in the CM category.  The splitting
results below are unconditional and do not use either zero theorem.
Only the assertions about the zeros of their CM factors use the
conditional input \textbf{M} of ZETA\_TWO CM01.  Its scope is all unitary
Hecke characters, including infinite-order characters; the source and
dependency audit of that input are carried from CM01 and MDOC01.
For a chosen finite-type integral model $\mathcal X$, the superscript
$S$ means the Euler product over $p\notin S$.  The full model zeta
function is always
\[
 \zeta_{\mathcal X}(s)=
       \left(\prod_{p\in S}Z(\mathcal X_p,p^{-s})\right)
                                  \zeta_X^S(s).
\]
The actual exceptional fibres therefore remain explicit; none is
replaced by a generic-fibre inertia factor without the correction in
ZETA\_TWO ZT04.

\subsection{NG01: a canonical maximal-CM summand, with descent}
\label{NG01}

\paragraph{Finding NG01.}
Let $M$ be a pure motive of abelian type over $\mathbb Q$, for absolute
Hodge cycles.  There is a canonical decomposition over $\mathbb Q$
\[
                    M=C(M)\oplus N(M),                         \tag{NG01.1}
\]
where $C(M)$ is potentially CM and $N(M)$ has no nonzero potentially-CM
motivic subobject.  Let $G_M$ be the algebraic group of the neutral
Tannakian category generated by $M$, for the Betti fibre functor, and put
$D_M=(G_M^\circ)^{\rm der}$.  Then
\[
 H_B(C(M))=H_B(M)^{D_M},\qquad
 H_B(N(M))=\langle (d-1)v:d\in D_M,\ v\in H_B(M)\rangle.       \tag{NG01.2}
\]
The second expression is interpreted after scalar extension to an
algebraic closure and then descended.  Equivalently it is the sum of the
nontrivial $D_M$-isotypic summands.  These formulas retain the action of
the entire, possibly disconnected group $G_M$.

Every $\ell$-adic realization of (NG01.1) is an actual direct sum of
Galois representations.  When $M=h^j(A)$, its CM part is therefore
available for the Hecke factorization of ZETA\_TWO CM01, and the quotient
of the good-prime local polynomial of $h^j(A)$ by that of $C(M)$ is the
local polynomial of $N(M)$.  This gives an exact factorization of the
incomplete $L$-function.

\paragraph{Proof.}
The category is semisimple neutral Tannakian by Deligne--Milne,
Proposition~6.5 \cite{NGDeligneMilne}.  The fully faithful Hodge realization over an
algebraically closed field (Theorem~6.25) and the descent sequence
(Proposition~6.23) identify $G_M^\circ$ with the Mumford--Tate group of
$H_B(M)$; see also Commelin, Remark~4.14(1) \cite{NGCommelin}.  In particular $G_M^\circ$
is reductive.  Its derived group $D_M$ is connected semisimple and is
characteristic in $G_M^\circ$, hence normal in $G_M$.

Restrict the representation $V=H_B(M)$ to $D_M$.  Complete reducibility
gives the direct sum of its trivial isotypic component $V^{D_M}$ and
the sum $V_{\ne1}$ of all its nontrivial isotypic components.  Normality
of $D_M$ implies that every $g\in G_M$ preserves the trivial isotypic
component and permutes the nontrivial ones.  Thus both summands and the
projection $e:V\to V^{D_M}$ are $G_M$-equivariant.  Tannaka duality
identifies $e$ with an idempotent endomorphism of $M$ over $\mathbb Q$.
This proves the descent of the projector, rather than merely its
existence after an unspecified finite extension.  Its image and kernel
give (NG01.1), and semisimplicity gives (NG01.2).

The identity component of the group acting on $V^{D_M}$ is a quotient
of the torus $G_M^\circ/D_M$.  An abelian motive with torus identity
component is in the category generated by Artin motives and potentially
CM abelian motives (Commelin, \S4.16, following Lemma~3.8 and
Theorem~4.13).  Hence $C(M)$ is potentially CM.  Conversely the connected
group acting on a potentially-CM subobject is a torus.  A morphism from
the connected semisimple group $D_M$ to this torus is trivial.  Such a
subobject is therefore contained in $V^{D_M}$, which proves maximality
and the assertion about $N(M)$.  The same argument proves uniqueness
and functoriality for morphisms: apply it in the common category
generated by source and target, where a morphism respects the trivial
and nontrivial isotypic components.  The formula using spans agrees with
$V_{\ne1}$ because its quotient is $D_M$-trivial and the category of
$D_M$-representations is semisimple.

A realization functor preserves direct sums and idempotents, so the
$\ell$-adic splitting follows.  At a good prime the determinant of
Frobenius on a direct sum is the product of the two determinants.  The
CM factor has the compatible Hecke description already used in
ZETA\_TWO CM01.  Division of the known good-prime polynomial of $h^j(A)$
by that polynomial is therefore a polynomial, independent of the
auxiliary $\ell$, and is the polynomial on $N(M)_\ell$.  This last
argument avoids assuming a separate compatibility theorem for every
possible absolute-Hodge projector.  Multiplication of the local
identities proves the asserted incomplete Euler-product identity.
\hfill$\square$

\paragraph{The precise role of Faltings.}
Faltings, \S5, S\"atze~3--4, printed p.~360 \cite{NGFaltings}, proves semisimplicity of
$H^1_\ell(A)$ and the endomorphism theorem.  Since the Zariski closure
$G_\ell$ acts faithfully on this semisimple representation, its
unipotent radical is trivial; in characteristic zero all its tensor
and exterior representations are semisimple.  Consequently
$H^j_\ell(A)=\bigwedge^jH^1_\ell(A)$ is semisimple for every $j$.
For clarity, the reductivity implication can be checked directly:
on every irreducible constituent the normal connected unipotent
radical has a nonzero fixed subspace, which is stable under $G_\ell$,
and therefore is the whole constituent.  Faithfulness forces that
radical to be trivial.  The finite component group causes no
semisimplicity problem in characteristic zero.

At a fixed $\ell$ there is also a canonical maximal potentially
abelian subrepresentation
\[
 \bigl(H^j_\ell(A)\bigr)^{(G_\ell^\circ)^{\rm der}}.           \tag{NG01.3}
\]
One must not identify (NG01.3) with $C(h^j(A))_\ell$ by citing
Faltings alone.  Deligne gives the inclusion
$G_\ell^\circ\subseteq\operatorname{MT}(A)_{\mathbb Q_\ell}$,
and hence inclusion of the motivic CM part in (NG01.3).
Equality follows, for example, when the Mumford--Tate conjecture is
known for $A$.  Faltings' endomorphism theorem for $H^1$ does not assert
that every projector in a higher tensor power is an absolute-Hodge
projector.  The proof of NG01 uses the motivic group and supplies the
required motivic splitting without making this identification.

\paragraph{Scope of the closure.}
NG01 closes the decomposition gap in ZETA\_TWO ZT06 whenever the pure
factor is of abelian type.  Deligne--Milne, Proposition~6.26 and
Corollary~6.27,
include curves, products with abelian and zero-dimensional varieties,
Fermat hypersurfaces, K3 surfaces, and the stated unirational varieties
of dimension at most three.  The conclusion is about motives for
absolute Hodge cycles.  It is not a general construction of Chow
projectors, and it does not assert such a splitting for an arbitrary
smooth projective variety outside the abelian-type category.

\subsection{NG02: the exact list for every curve}
\label{NG02}

\paragraph{Finding NG02.}
Let $C/\mathbb Q$ be a reduced separated finite-type curve.  Zero-dimensional
components can be added as separate Artin terms, so first assume its
components have dimension one.  Passing to the reduction does not
change the zeta function, so this also treats nonreduced curves.
Let $\widetilde C\to C$ be its
normalization, $\overline C$ the smooth projective compactification of
$\widetilde C$, $D=\overline C\setminus\widetilde C$, $Z$ the finite
singular locus of $C$, and $\widetilde Z$ its inverse image.  All finite
schemes in the following formula are reduced.  Put
\[
 \Pi=H^0(\overline C_{\overline{\mathbb Q}},\mathbb Q),\qquad
 B=H^0(Z)-H^0(\widetilde Z)-H^0(D),\qquad
 J=\operatorname{Pic}^0(\overline C/\mathbb Q).
\]
Thus $\Pi$ is the permutation Artin motive of geometric connected
components and $B$ is a virtual Artin motive.  For a sufficiently
large finite $S$, the incomplete zeta function of a spread-out model is
\[
 \zeta^S_C(s)=
 \frac{L^S(\Pi+B,s)L^S(\Pi,s-1)}{L^S(H^1(J),s)}.             \tag{NG02.1}
\]
Write the $\mathbb Q$-isogeny decomposition
\[
                     J\sim_{\mathbb Q}\prod_i B_i^{m_i},
\]
with $B_i$ pairwise nonisogenous $\mathbb Q$-simple abelian varieties,
and let $I_{\rm nc}$ consist of those that are not geometrically of CM
type.  The complete non-CM list in (NG02.1) is precisely
\[
            \bigl\{\,L^S(B_i,s)\text{ with multiplicity }m_i:
                                i\in I_{\rm nc}\,\bigr\}.     \tag{NG02.2}
\]
Each factor has weight one and its desired RH line is $\Re s=1$.
All corrections for singularities, missing points, and geometric
components are already in the explicitly displayed Artin factors.

\paragraph{Proof.}
Normalization is an isomorphism away from the indicated finite closed
sets.  The scissor relation gives
\[
 [C]=[\overline C]-[D]-[\widetilde Z]+[Z]
                  \quad\text{in }K_0(\operatorname{Var}_{\mathbb Q}).
\]
The relations spread over a sufficiently small open of
$\operatorname{Spec}\mathbb Z$ by the same finite-presentation argument
used in ZETA\_TWO ZT06 (EGA IV$_3$, (8.8.2)).  For a smooth projective
curve the pure cohomology is $\Pi$, $H^1(J)$, and $\Pi(-1)$ in degrees
zero, one, and two.  The geometrically connected case is
Deligne--Milne, Proposition~6.14; for general components the same
identity descends with the permutation action, or equivalently follows
by restriction of scalars from their fields of definition.  No
rational base point is needed: changing a geometric Abel--Jacobi map
by translation does not change its map on $H^1$.  Applying
the trace formula to the spread-out scissor relation proves (NG02.1).

The isogeny decomposition gives the direct sum of the $H^1(B_i)$ with
the stated multiplicities.  The sum of all isotypic factors that
become CM is stable under Galois, so its central idempotent is defined
over $\mathbb Q$.  Equivalently it is the degree-one projector of NG01.
Deligne--Milne, Proposition~6.21, identifies degree-one motives with
abelian varieties up to isogeny; therefore a nonzero CM submotive of
$H^1(B_i)$ would give a nonzero CM isogeny factor of $B_i$.  Simplicity
forces the whole $B_i$ to have potential CM.  Thus no such subobject
exists for $i\in I_{\rm nc}$, and every other factor belongs to CM01.
This proves completeness of (NG02.2).  Faltings' theorem verifies the
corresponding degree-one semisimple Galois decomposition, including
the endomorphism idempotents.  No splitting of the mixed
compact-support cohomology of an open or singular curve has been
assumed; its Grothendieck class is what the scissor calculation uses.
\hfill$\square$

\paragraph{Extension to nonseparated curves.}
The same exact non-CM list applies to any finite-type
$\mathbb Q$-scheme $C$ of dimension at most one, without a
separatedness assumption.  Reduce $C$ first.  For every
one-dimensional irreducible component choose a dense affine open
$U_i$ contained in that component and avoiding all other components.
These opens are also open in $C$ and are pairwise disjoint: take
an affine neighbourhood of its generic point inside the open
complement of the other components.  Put $U=\coprod_iU_i$ and
$Z_0=(C\setminus U)_{\rm red}$.  Every one-dimensional generic point
lies in $U$, so $Z_0$ is zero-dimensional and finite over $\mathbb Q$.
For this last assertion no separatedness is needed: a zero-dimensional
finite-type scheme has finitely many open-and-closed points, each
with an Artinian finite-dimensional coordinate algebra over the base.
Thus $U$ is separated, and the closed-point scissor identity gives,
after enlarging $S$,
\[
               \zeta_C^S(s)=\zeta_U^S(s)L^S(H^0(Z_0),s).
                                                               \tag{NG02.4}
\]
Apply (NG02.1) to $U$: its virtual Artin term $B_U$ is simply
replaced by $B_U+H^0(Z_0)$.  The residual list (NG02.2) and the
order identity (NG02.3) are unchanged with this replacement.
Indeed the added factor is Artin, and the smooth projective
normalizations of the $U_i$ depend only on the function fields of
the original components.  Consequently the Jacobian factors in
the list do not depend on the affine opens chosen.  If $U$ is empty,
the entire zeta function is Artin.  This proves the extension by
separated strata and does not require a proper compactification
of a nonseparated normalization.\hfill$\square$

\paragraph{A precise coefficient refinement.}
For a $\mathbb Q$-simple $B_i$, let $E_i$ be the centre of the division
algebra $\operatorname{End}^0_{\mathbb Q}(B_i)$ and write its dimension
over $E_i$ as $d_i^2$.  Over $\overline{\mathbb Q}_\ell$, Faltings'
endomorphism theorem and the double-centralizer decomposition give
$[E_i:\mathbb Q]$ pairwise nonisomorphic irreducible constituents, each
of rank
\[
                  r_i=\frac{2\dim B_i}{d_i[E_i:\mathbb Q]},
\]
and each of multiplicity $d_i$.  Indeed the rational Betti realization
is a free module over the division algebra, so its dimensions at all
embeddings of $E_i$ are equal.  The commutant becomes the
product of $[E_i:\mathbb Q]$ matrix algebras
$M_{d_i}(\overline{\mathbb Q}_\ell)$; comparison with the commutant of
a semisimple representation proves both assertions.  This identifies
the degree-two (GL$_2$-type) systems and the higher-rank systems in
(NG02.2) without calling every abelian-variety factor a GL$_2$
$L$-function.  If a coefficient refinement or automorphic realization
is known, it may be used; the rational $L(B_i,s)$ is already an exact
named factor and needs no such additional identification.

\paragraph{Noncancellation for the curve product.}
Under \textbf{M}, and assuming meromorphic continuation of the displayed
non-CM product, every point with
$\tfrac12<\Re s<\tfrac32$ and $\Re s\ne1$ satisfies
\[
 \operatorname{ord}_s\zeta^S_C
       =-\operatorname{ord}_s
             \prod_{i\in I_{\rm nc}}L^S(B_i,s)^{m_i}.          \tag{NG02.3}
\]
To check this, the weight-zero Artin factors have their nontrivial
divisors on the boundary $\Re s=1/2$, and their weight-two twists on
the boundary $\Re s=3/2$.  Their possible real endpoint divisors in
this open strip are on $\Re s=1$.  The CM weight-one factors have
their global divisors on $\Re s=1$ by CM01; their omitted local
Hecke factors have weight one and therefore divisors on
$\Re s=1/2$.  Deleted factors of the Artin terms lie on real parts
zero or one.  Thus all other terms are invertible at the indicated
points and the order identity follows directly from (NG02.1).
This version with incomplete functions needs no new assertion about
the complete bad-prime factors of a general non-CM abelian variety.
When individual non-CM factors are known to be holomorphic
in the strip, their positive multiplicities prevent cancellation
among them.  Without holomorphy, the exact converse is for the
aggregate meromorphic product, not for every separate factor.

\subsection{NG03: the exact abelian-variety list}
\label{NG03}

\paragraph{Finding NG03.}
Let $A/\mathbb Q$ have dimension $g$ and choose a polarization over
$\mathbb Q$.  Let $P^r(A)$ denote primitive degree-$r$ cohomology,
$0\le r\le g$, and set
\[
                  C_r=C(P^r(A)),\qquad N_r=N(P^r(A)).
\]
These are absolute-Hodge motives over $\mathbb Q$.  Then
\[
 \zeta_A^S(s)=
 \prod_{r=0}^g\prod_{k=0}^{g-r}
 \left\{L^S(C_r,s-k)L^S(N_r,s-k)\right\}^{(-1)^r}.            \tag{NG03.1}
\]
Here $P^0=C_0=\mathbb Q(0)$ and $N_0=0$.  Consequently the exact
residual list after the CM part of ZT06 is removed is
\[
 \boxed{\ \left\{L^S(N_r,s-k):1\le r\le g,\quad
                    0\le k\le g-r,\quad N_r\ne0\right\}.\ } \tag{NG03.2}
\]
The exponent is $(-1)^r$, the weight is $r+2k$, and the desired line is
$\Re s=(r+1)/2+k$.  The remaining analytic requirements are precisely
continuation, functional equation and divisor RH for these factors,
to the extent not known independently, together with the explicit
finite local factors of ZT04.  Under \textbf{M} the CM factors satisfy
the conclusion of CM01; no assertion that \textbf{M} settles
(NG03.2) is made.

\paragraph{Proof.}
The cohomology algebra is the exterior algebra of $H^1(A)$
(Deligne--Milne, Proposition~6.13).  The symplectic Lefschetz
decomposition is
\[
 h^j(A)=\bigoplus_{\substack{r+2k=j\\0\le r\le g\\0\le k\le g-r}}
                         P^r(A)(-k).                         \tag{NG03.3}
\]
One may construct it directly on the exterior algebra: in a symplectic
basis, wedge multiplication by the polarization and contraction by
the inverse alternating form, with the degree operator, form the
standard $\mathfrak{sl}_2$ triple.  Every irreducible
$\mathfrak{sl}_2$-string beginning in degree $r\le g$ has the degrees
$r,r+2,\ldots,2g-r$.  Its lowest-degree vectors are exactly the kernel
of contraction, namely $P^r$; successive multiplication by the
polarization supplies the indicated Tate twists.  Complete
reducibility gives (NG03.3).  The operators are Hodge morphisms and
are absolute Hodge by Deligne's theorem, so the decomposition exists
in the chosen motivic category over $\mathbb Q$.

Apply NG01 to each primitive object.  Tate twists do not change
whether the connected derived group acts trivially, hence
$C(P^r(-k))=C_r(-k)$ and $N(P^r(-k))=N_r(-k)$.  Frobenius determinants
multiply on these direct sums.  Degree $j=r+2k$ has parity $r$, and
$L(M(-k),s)=L(M,s-k)$.  The trace-formula Euler product now gives
(NG03.1).  The list (NG03.2) and its weight and signs follow.  As in
ZT06, this is an exact factorization and a sufficient analytic gate;
it does not claim a converse that would rule out meromorphic
cancellation between distinct residual factors.\hfill$\square$

\paragraph{What is, and is not, computed uniformly.}
For a specified Mumford--Tate representation $V=H^1_B(A)$, (NG03.2)
is finite linear algebra.  If $X_1,\ldots,X_t$ span
$\operatorname{Lie}(\operatorname{MT}(A)^{\rm der})$, then
\[
             H_B(C_r)=\bigcap_{a=1}^t
                  \ker\bigl(d\rho_r(X_a):P^r_B\to P^r_B\bigr).
                                                                    \tag{NG03.4}
\]
The derived group is connected, so Lie-algebra invariants and group
invariants agree.  The complementary isotypic subspace then gives
$N_r$, and local determinants give its Euler factors.  No general
effective algorithm recovering the full Mumford--Tate representation
of an arbitrary abelian variety from equations is proved by
Faltings or Deligne.  Thus (NG03.2)--(NG03.4) are an exact universal
formula with a specified input, not a claimed algorithm solving that
separate identification problem.  If only $A\sim A_{\rm CM}\times
A_{\rm nc}$ is known, the correct formula is
\[
 C(h^j(A))=
 \bigoplus_{a+b=j}\bigwedge^a h^1(A_{\rm CM})\otimes
                       C\!\left(\bigwedge^b h^1(A_{\rm nc})\right).
                                                                    \tag{NG03.5}
\]
Indeed the derived group acts trivially on the first factor and
surjects onto the derived group acting on the second.  This formula
does not discard the CM summands created inside exterior powers of
the non-CM part.

\subsection{NG04: the full symplectic case}
\label{NG04}

\paragraph{Finding NG04.}
Suppose $\operatorname{MT}(A)=\operatorname{GSp}_{2g}$ in its standard
representation.  Then $C(h^{2i}(A))=\mathbb Q(-i)$,
$C(h^{2i+1}(A))=0$, and $N_r=P^r$ for $1\le r\le g$.
In particular
\[
 \zeta_A^S(s)=\prod_{i=0}^g\zeta^S(s-i)
       \prod_{r=1}^g\prod_{k=0}^{g-r}L^S(P^r(A),s-k)^{(-1)^r},
 \qquad \operatorname{rank}P^r={2g\choose r}-{2g\choose r-2}. \tag{NG04.1}
\]
The non-CM functions are the primitive exterior-power Hasse--Weil
$L$-functions and their indicated shifts.  For a full-symplectic
abelian surface this means $L(H^1(A),s)$,
$L(H^1(A),s-1)$, and the rank-five primitive exterior-square function
$L(P^2(A),s)$; the lines are $1$, $2$, and $3/2$, respectively.

\paragraph{Proof.}
The derived group is $\operatorname{Sp}_{2g}$.  We give a direct check
of its invariants in the exterior algebra, so no assumption about
unknown algebraic cycles is involved.  Choose a symplectic basis
$e_1,f_1,\ldots,e_g,f_g$.  The diagonal torus sends $e_i$ to $t_i e_i$
and $f_i$ to $t_i^{-1}f_i$.  A zero-weight exterior monomial must
contain both $e_i,f_i$ or neither, for every $i$; hence no odd-degree
invariant exists.  In degree $2i$, the zero-weight monomials are the
products of $i$ distinct pairs $e_a\wedge f_a$.  The subgroup
permuting the symplectic pairs acts transitively on these monomials,
with positive sign because each pair has degree two.  Thus any
$\operatorname{Sp}_{2g}$-invariant is a scalar multiple of their sum.
That sum is $1/i!$ times the $i$-th power of the polarization bivector,
which is indeed invariant.  The invariant space therefore has
dimension exactly one.  NG01 identifies it with $\mathbb Q(-i)$,
using the polarization defined over $\mathbb Q$.

For $r\le g$, wedge multiplication by the polarization is injective
from degree $r-2$ to degree $r$.  Its image and the primitive kernel
form a direct sum by the symplectic Lefschetz decomposition in NG03.
This gives the displayed rank.  In positive even primitive degree
the sole invariant, the power of the polarization, is in the image
of wedge multiplication and hence has zero intersection with the
primitive kernel.  In odd degree there is no invariant at all.
Consequently $C(P^r)=0$ for every $r\ge1$.  Substitute in NG03.
\hfill$\square$

\subsection{NG05: every power of a non-CM elliptic curve}
\label{NG05}

\paragraph{Finding NG05.}
Let $E/\mathbb Q$ be non-CM and $m\ge1$.  Adopt the convention that
$\binom m t=0$ outside $0\le t\le m$.  For $0\le b\le a\le m$, put
\[
 d_m(a,b)=\binom ma\binom mb-
                        \binom m{a+1}\binom m{b-1}.           \tag{NG05.1}
\]
The complete decomposition is
\[
 h^j(E^m)=\bigoplus_{\substack{0\le b\le a\le m\\a+b=j}}
       \left(\operatorname{Sym}^{a-b}h^1(E)(-b)\right)^{\oplus d_m(a,b)}.
                                                                    \tag{NG05.2}
\]
The CM terms are exactly $a=b$, and all other terms are non-CM.
Consequently
\[
 \zeta_{E^m}^S(s)=
 \prod_{0\le b\le a\le m}
 L^S(\operatorname{Sym}^{a-b}E,s-b)^{(-1)^{a+b}d_m(a,b)},       \tag{NG05.3}
\]
where $L(\operatorname{Sym}^0E,s)=\zeta(s)$.  The distinct non-CM
analytic inputs, before their specified shifts, are exactly
$L(\operatorname{Sym}^nE,s)$, $1\le n\le m$.  Their motivic RH lines
are $\Re s=(n+1)/2$.  Their analytic continuation is already known:
for $n\ge2$ see Newton--Thorne, \emph{Symmetric power functoriality
for holomorphic modular forms, II}, Theorem~A and Corollary~B
(printed pp.~117--118) \cite{NGNewtonThorne}; the $n=1$ case is the modularity theorem
used in ZETA\_TWO ZT07.  Their RH is an additional analytic assertion.

\paragraph{Proof.}
Write $V=h^1(E)$ and $W=\mathbb Q^m$ with trivial action.
Then $h^j(E^m)=\bigwedge^j(V\otimes W)$.  The Mumford--Tate group
of a non-CM elliptic curve is $\operatorname{GL}_2$ (for example,
Moonen, Example~2.1.13 \cite{NGMoonen}); its derived group is $\operatorname{SL}_2$.
We compute the decomposition directly.  On a diagonal matrix with
eigenvalues $x,y$, the exterior-character generating series is
\[
                      (1+tx)^m(1+ty)^m.
\]
The coefficient of $t^{a+b}x^ay^b$ is
$w(a,b)=\binom ma\binom mb$.  The irreducible
$\operatorname{GL}_2$-representation of highest weight $(a,b)$,
$a\ge b$, has character
$x^ay^b+x^{a-1}y^{b+1}+\cdots+x^by^a$ and is
$\operatorname{Sym}^{a-b}V\otimes(\det V)^b$.
For fixed total degree, subtraction of the adjacent weight
multiplicities $w(a,b)-w(a+1,b-1)$ therefore isolates its highest-weight
multiplicity.  This is (NG05.1).  The exterior representation is
semisimple, so this character calculation gives its direct-sum
decomposition.  Equivalently, it is the two-row case of the exterior
Cauchy decomposition.  The corresponding Hodge-equivariant
projectors are absolute Hodge and descend as morphisms in the
$\mathbb Q$-category generated by $V$: its full motivic group is a
subgroup of $\operatorname{GL}(V)$ whose identity component is already
$\operatorname{GL}(V)$, so it equals that group.  Thus the identity is motivic.
Since $\det V=\mathbb Q(-1)$, it is exactly (NG05.2).

For $a>b$, $\operatorname{Sym}^{a-b}V$ is a nontrivial irreducible
$\operatorname{SL}_2$-representation, hence has no derived-group
invariants.  For $a=b$ it is the trivial representation, followed by
a Tate twist.  NG01 proves the asserted complete CM/non-CM
classification.  Taking determinants in degree $j=a+b$ gives
(NG05.3).\hfill$\square$

In particular the first example beyond a curve is
\[
 \zeta_{E^2}^S(s)=
 \frac{\zeta^S(s)\,\zeta^S(s-1)^3\,\zeta^S(s-2)
                        L^S(\operatorname{Sym}^2E,s)}
      {L^S(E,s)^2 L^S(E,s-1)^2}.                              \tag{NG05.4}
\]
Thus the non-CM gate for $E^2$ contains a symmetric-square
$L$-function, even though the only simple abelian isogeny factor is
$E$.  Conversely its degree-two cohomology has three Tate summands,
even though $h^1(E^2)$ has no CM summand.  This refutes the procedure
of removing only the exterior algebra of the CM isogeny factors and
calling everything left non-CM.

\subsection{NG06: non-Tate CM terms can occur in higher degree}
\label{NG06}

\paragraph{Finding NG06.}
Let $A/\mathbb Q$ have dimension $g$, and suppose that over a number
field $F$ it has an action by an imaginary quadratic field $K$.
Suppose that the $K$-action on $H^1(A)$ has signature $(p,q)$, where
$p+q=g$.  Over $F$, the rational Hodge subspace
\[
                         W=\bigwedge_K^g H^1(A)               \tag{NG06.1}
\]
has rank two over $\mathbb Q$, occurs as a summand of $H^g(A)$,
and has Hodge types $(p,q)$ and $(q,p)$.  It is a CM absolute-Hodge
motive.  If $p\ne q$, it is not Artin--Tate.  Its Galois-stable sum of
conjugates in $h^g(A)_{\overline{\mathbb Q}}$ descends to a CM summand
over $\mathbb Q$, and is included by NG01.

\paragraph{Proof.}
Over a splitting field for $K$, write
$H^1=V_\sigma\oplus V_{\bar\sigma}$, with each summand of dimension
$g$.  The exterior algebra decomposes by its bidegree in these two
spaces.  The sum
\[
                      \bigwedge^g V_\sigma\ \oplus\
                      \bigwedge^g V_{\bar\sigma}
\]
is invariant under the coefficient Galois group and defines the
rational direct summand (NG06.1).  Taking determinants adds the Hodge
indices, so its two types are precisely $(p,q)$ and $(q,p)$.
The projector onto these two extreme exterior degrees, along all
the mixed degrees, is invariant under $G_F$ because the $K$-action
is defined over $F$.  Its absolute-Hodge realization thus descends
over $F$.
The summand is a free rank-one $K$-module.  Its Mumford--Tate group
commutes with $K$ and is consequently contained in
$\operatorname{Res}_{K/\mathbb Q}\mathbb G_m$.  It is therefore a
torus.  Deligne's theorem makes its Hodge projector an absolute-Hodge
projector, and the CM characterization used in NG01 proves the CM
assertion.  An Artin--Tate pure Hodge structure has only diagonal
types, so $p\ne q$ excludes that possibility.

The projector is defined over a finite extension by Deligne,
Proposition~2.9 \cite{NGDeligne82}.  There are only finitely many conjugate subobjects.
Their sum is Galois stable, remains CM because the category is closed
under direct sums and quotients, and has a projector in the semisimple
category.  Conjugation preserves the CM category because it preserves
Artin motives and the endomorphism criterion for CM abelian varieties.
Choose a finite Galois extension defining the sum and a projector
onto it, and average the projector over its finite Galois group.
Every conjugate projector has image in the same subobject and is the
identity on it.  Their average consequently has the same image and
is still idempotent; it is Galois invariant and descends over
$\mathbb Q$.  This proves the final assertion without assuming the
original $K$-action is defined over $\mathbb Q$.\hfill$\square$

For example, signature $(2,1)$ in dimension three yields a weight-three
CM term of types $(2,1)$ and $(1,2)$.  The statement is a construction
for every abelian variety with the specified action and signature;
it does not claim that an arbitrary abelian threefold has that action.
Together with NG05, it explains why higher CM terms must be selected
degree by degree and cannot be read solely from a list of CM
degree-one factors.

\subsection{NG07: an explicit non-CM Jacobian with non-Tate CM cohomology}
\label{NG07}

\paragraph{Finding NG07.}
Let $C/\mathbb Q$ be the smooth projective curve with affine equation
\[
                            y^3=x^5-x-1,
\]
and let $A=\operatorname{Jac}(C)$.  Then $A$ is an absolutely simple
non-CM abelian fourfold and $h^1(A)$ has no nonzero CM summand.
Nevertheless $h^4(A)$ contains a rank-two CM summand $W$ defined over
$\mathbb Q$, with Hodge types $(3,1)$ and $(1,3)$.  In particular
$W$ is not Artin--Tate.  It occurs in the primitive part $P^4(A)$
for the canonical Jacobian polarization.

\paragraph{Proof.}
Zarhin, Theorem~1.1 and Examples~1.3(1) of
\emph{The endomorphism rings of jacobians of cyclic covers of the
projective line} \cite{NGZarhin}, apply to the exact family
$y^p=x^n-x-1$ with $n\ge5$ and $p$ odd.  For $n=5,p=3$ they give
\[
 \operatorname{End}_{\overline{\mathbb Q}}(A)=\mathbb Z[\zeta_3],
 \qquad \dim A=(3-1)(5-1)/2=4.
\]
In particular $A$ is absolutely simple.  Its geometric endomorphism
algebra has dimension two over $\mathbb Q$, so cannot contain a
commutative semisimple subalgebra of dimension $2\dim A=8$; thus
$A$ is non-CM.  NG02, or the degree-one equivalence in
Deligne--Milne 6.21, shows that $h^1(A)$ has no CM subobject.
For separate input checks, the discriminant of $x^5-x-1$ was computed as
$2869$, its reduction modulo $3$ is irreducible, and its reduction
modulo $2$ is
$(x^2+x+1)(x^3+x^2+1)$.  These computations are recorded separately;
the Galois-group conclusion used here is Zarhin's stated example.

Put $K=\mathbb Q(\zeta_3)$ and let
$\delta(x,y)=(x,\zeta_3y)$.  A basis of regular differentials is
\[
 \frac{dx}{y},\quad \frac{dx}{y^2},\quad
                    \frac{x\,dx}{y^2},\quad
                    \frac{x^2\,dx}{y^2}.                    \tag{NG07.1}
\]
For completeness, at a finite branch point $y$ is a local parameter
and $dx$ has order two, so all four forms are regular.  At the unique
point at infinity, $x,y,dx$ have orders $-3,-5,-4$, respectively.
The forms in (NG07.1) therefore have orders $1,6,3,0$, respectively,
and are regular there.  They are linearly independent: the first has
$\delta$-eigenvalue $\zeta_3^2$, while the remaining three have
eigenvalue $\zeta_3$ and are independent powers of $x$ times the same
nonzero differential.  The genus is four, so they form a basis.
This proves that the $K$-signature is $(3,1)$, up to exchanging the
two embeddings.

NG06 now supplies
$W=\bigwedge_K^4 H^1(A)$, of types $(3,1)$ and $(1,3)$.
Here $K$ is the entire geometric endomorphism algebra, so Galois acts
on it by field automorphisms, either fixing or exchanging its two
embeddings.  The sum of the two extreme exterior eigenspaces that
defines $W$ and its complementary mixed-degree subspace are
therefore stable under the full Galois group.  Their absolute-Hodge
projector is defined over $\mathbb Q$, so $W$ itself, with rank two,
descends; it is not necessary to enlarge it to a sum of conjugates.

Finally $\delta$ preserves the canonical Jacobian polarization.
Consequently the polarization pairs the two $K$-eigenspaces in
$H^1$ with one another and has no same-eigenspace component.  The
contraction operator therefore kills both
$\bigwedge^4 V_\sigma$ and $\bigwedge^4 V_{\bar\sigma}$.
This puts $W$ in $P^4(A)$.  Its non-diagonal Hodge types prove that
it is not Artin--Tate.\hfill$\square$

\paragraph{Placement in Grothendieck's frameworks.}
Over $U=\operatorname{Spec}\mathbb Z[1/S\ell]$, for separated models
the formulas are
identities in $K_0(D_c^b(U,\mathbb Q_\ell))$ for the realizations of
$Rf_!\mathbb Q_\ell$, with actual pure direct-sum projectors in the
abelian-type cases and a virtual scissor identity for open or singular
curves.  The nonseparated-curve extension (NG02.4) uses the finite sum
of the compact-support classes of its separated strata.
The support and trace framework is SGA~4 XVII,
\S\S5.1--5.3, and SGA~5 XV, \S3, no.~2, equation~(1), as located and
used in ZETA\_TWO ZT04--ZT06.  Exterior powers, determinants and
their passage to Grothendieck groups are the $\lambda$-ring mechanism
of SGA~6 V, \S2 and \S7, (7.1.1), and ZETA\_TWO ZT08.
The relevant documented archive locations are cote~60, scans~3 and~5
for determinant/exterior-power algebra, scan~7 for reciprocity, and
scan~9 for the explicit Weil and dual/Tate hypotheses.  Cote~34,
scans~64 and~74, remains the finite-field Riemann--Weil context
recorded in the predecessor.  These are carried archival locators,
not new readings of unavailable scans, and none is attributed the
later Faltings--Deligne splitting theorem.  In the runs, NG01--NG03
close the specific abelian-type decomposition gap noted in
INTAKE\_ZETA\_TWO, \S1, row ZT06 and \S5.1; NG05 refines the first
non-CM example ZT07; NG01 and NG06--NG07 extend the usable scope of CM01.


\clearpage
\section{Framework and predecessor locator register}

The following register states where the language and the predecessor
questions live.  It distinguishes a framework locator from a claim
that a new theorem occurs at that locator.

\small
\begin{longtable}{p{.19\textwidth}p{.37\textwidth}p{.36\textwidth}}
\toprule
Findings & Scheme, motive, or determinant interpretation & Exact supplied or classical anchor\\
\midrule
\endhead
NX01--NX04; NC01 &
$\zeta_{\operatorname{Spec}\mathcal O_K}$, the finite Picard group,
and the logarithmic unit lattice; explicit size bounds, not a recovery
of integral $K$-theory from zeros. &
ZETA\_TWO IA05, ZT08, SV01--SV02; SA06 B.4a--e in its intake
\S4.1. No cote/scan for B.4a--e is supplied.\\[3pt]
NC02--NC04 &
The metrized determinant of the Hodge bundle on the abelian moduli
problem; isogenies of CM motives and indices of central integral
orders. &
ZETA\_TWO IA05 and intake \S6(1)(a);
SGA~6 V, \S2 and (7.1.1), for the determinant framework only.
No EGA/SGA/FGA theorem is asserted to be the new numerical bound.\\[3pt]
Quadratic specialization &
Class numbers and the existing supersingular-count input. &
SA06 A.8a; cote~15, scan~70 as carried in the ZETA\_TWO intake
\S5.4; IA02--IA03. A.8c is a related row whose ring $R$ is not
defined by the supplied intake.\\[3pt]
NX-L01--NX-L07 &
Frobenius lifts on flat schemes; actual special-fibre
$\ell$-adic cohomology, compact supports, and local determinants. &
Borger \S1.1, Theorem~6.1, Corollary~6.2;
SGA~4 XII 5.1--5.2, XVI 2.1, XVII \S\S5.1--5.3;
SGA~5 XV \S3 no.~2 equation~(1);
ZETA\_TWO ZT04--ZT05.\\[3pt]
NX-L01; binomial scope &
Commuting Frobenius lifts versus the special identity-Adams
binomial structure. &
DP\_MULT DP11; its intake, verdict table.
DP\_MULT reports cote~153, PDF scan~3/archive p.~2
for binomial polynomials.\\[3pt]
NG01--NG07 &
Absolute-Hodge motives, canonical idempotents, exterior powers,
Lefschetz strings, and compatible local determinants. &
Deligne--Milne 6.5, 6.13--6.14, 6.21, 6.23, 6.25--6.27;
Faltings \S5 S\"atze~3--4; ZETA\_TWO ZT06--ZT08, CM01.\\[3pt]
NG02 and models &
Scissor classes spread to a good open; singular and open curves
give virtual Artin corrections. &
EGA IV$_3$ (8.8.2); SGA~4 XVII \S\S5.1--5.3;
ZETA\_TWO ZT06.\\[3pt]
All determinant translations &
Exterior powers, reciprocal characteristic polynomials, and the
explicit separation of purity from determinant algebra. &
Cote~60 scans~3 and~5; scan~7 (archive p.~6);
scan~9 (archive p.~8), the latter explicitly assuming Weil
and dual/Tate data. ZETA\_TWO AR01, ZT08.\\[3pt]
Weight terminology &
Finite-field Riemann--Weil statements are distinct from global
non-CM RH. &
Cote~34 scan~64, struck Remark~1.17;
scan~74, (2.5.5). Readings carried from ZETA\_TWO AR01
and its intake \S\S1,5.5; no transcription change.\\
\bottomrule
\end{longtable}
\normalsize

\part{Unchecked findings and boundaries}
\section{Check status and the residual analytic questions}

\paragraph{Unchecked findings.}
There are no unchecked findings included in this paper.
Each proof in Part~I has a recorded whole-proof reading in
\path{CHECK_EVENTS.tsv}; important cross-checks have a second reader.
These are targeted mathematical checks, not proof-assistant
formalizations.  The finite computations are auxiliary certificates
of the specific arithmetic or numerical assertions they check.

\paragraph{What remains outside the proved assertions.}
The following limits are part of the exact conclusions, not silently
assumed missing lemmas.
\begin{enumerate}
\item The new residue estimate is a zero-exclusion and special-value
statement.  It does not prove the full RH of a nonabelian
Dedekind zeta function.
\item The class-number and CM-orbit constants are explicit.
The general o-minimal counting and family-reduction statements
used elsewhere in Tsimerman's proof do not become effective algorithms
merely by substituting this orbit bound.  The cited modular-curve
effective replacement is distinguished in the CM section.
\item The literal $\Lambda$-scheme conjecture is refuted, and the
proper-flat question is proved.  A new version restricted to
separated nonproper schemes is a different question: the present
gluing example does not address it, and the proper tower proof
does not assert a $\Lambda$-equivariant compactification.
\item The abelian-type CM/non-CM splitting is motivic, but it is
not a claim that every higher-tensor $\ell$-adic projector is a
motivic projector.  NG01 supplies the required splitting without
assuming the Mumford--Tate conjecture.
\item NG02--NG05 name the non-CM analytic inputs exactly.
Their divisor RH, and in general the continuation and functional
equations not already known, remain the analytic gate.
For powers of a non-CM elliptic curve the cited symmetric-power
functoriality removes the continuation gate.
\item The CM-factor RH conclusions in the non-CM section are
conditional on M in its stated infinite-order scope.
The supplied draft's unnumbered claim is not promoted to a completed
numbered theorem.
\end{enumerate}

The task's three-hour stopping review was not required to decide any
of the three requested mathematical outputs: they reached the
proved/refuted/reduction states recorded in the ledger.
The ledger records actual elapsed-time checkpoint and review events,
without treating document assembly as mathematical movement.

\clearpage
\begin{thebibliography}{99}
\addcontentsline{toc}{section}{Bibliography and exact source locators}


\bibitem{Alpoge}
L.~Alp\"oge, \emph{Remark}, supplied four-page version of
20 November 2026, step~20, p.~4.  Endpoint declarations in
\path{ALPOGE_Main.lean}, with their predicate definitions in
\path{ALPOGE_Statements.lean}, as specified in the introduction.
\url{https://alpo.ge/remark.pdf}.  The supplied theorem status and
ZETA\_TWO IA01 audit are accepted as instructed.

\bibitem{Masullo}
A.~M.~Masullo, \emph{The Music of the Primes}, supplied 940-page
draft, 20 November 2026.  The claim used is \S1.2, PDF p.~12,
unnumbered.  The intended but empty divisor endpoint is
Theorem~13.3.1.2, PDF p.~522; compare \S14.1.2, PDF p.~528.
The supplied \path{MASULLO_STATEMENTS_INDEX.md} and ZETA\_TWO
MDOC01 fix this documentary status.

\bibitem{Anderson1986}
G.~W.~Anderson, \emph{Cyclotomy and an extension of the Taniyama
group}, Compositio Math. \textbf{57} (1986), 153--217.
\S5.3, p.~178; Corollary~5.7.2 and Remark~5.7.3, p.~181.
\url{https://www.numdam.org/item/CM_1986__57_2_153_0.pdf}.
The same-weight factorization is carried through ZETA\_TWO CM01.

\bibitem{ZetaTwo}
\emph{ZETA\_TWO}, supplied return
\path{RETURN_ZETA_TWO_final_20261121T0032Z/}.
Paper statements ZT01, ZT04--ZT08, CM01, IA01--IA05,
SV01--SV02, AR01 and MDOC01; the existing
\path{CONDITIONAL_AUDIT.tsv}; and
\path{INTAKE_ZETA_TWO_20261121.md},
\S1 (verdicts), \S4 (framework), \S5.1 (ZT06 gap),
\S5.4 (SA06), \S5.5 (archive readings), \S6(1)(a--c)
(the present questions).

\bibitem{DPMult}
\emph{DP\_MULT}, supplied
\path{RETURN_DP_MULT_checkpoint_03_20261121T0205Z/},
statement DP11; \path{INTAKE_DP_MULT_20261121.md},
\S1 verdict for DP11 and \S5 programme links.
DP11 is the binomial/identity-Adams criterion, not a claim that
every torsion-free $\lambda$-ring has identity Adams operations.

\bibitem{SA06}
\path{INTAKE_RESEARCH_SA06_20261120.md}, \S1 rows
A.8a--c and B.4a--g; \S4.1, final class-number paragraph,
and \S4.2 rows B.4a--e.  In particular B.4e supplies the
sufficient-power expression carried here, while B.4c remains
refuted.  No exact manuscript cote and scan for B.4a--e
is supplied in this intake.

\bibitem{SGA4}
M.~Artin, A.~Grothendieck, J.-L.~Verdier et al.,
\emph{Th\'eorie des topos et cohomologie \'etale des sch\'emas},
SGA~4. Tome~3, Lecture Notes in Math.~305 (1973):
Expos\'e XII, Th\'eor\`eme~5.1 and Corollaire~5.2;
Expos\'e XVI, Th\'eor\`eme~2.1;
Expos\'e XVII, \S\S5.1--5.3.

\bibitem{SGA5}
A.~Grothendieck et al., \emph{Cohomologie $\ell$-adique et
fonctions $L$}, SGA~5, Lecture Notes in Math.~589 (1977),
Expos\'e XV, \S3, no.~2, unnumbered Th\'eor\`eme, equation~(1).

\bibitem{SGA6}
P.~Berthelot, A.~Grothendieck, L.~Illusie et al.,
\emph{Th\'eorie des intersections et th\'eor\`eme de Riemann--Roch},
SGA~6, Lecture Notes in Math.~225 (1971),
Expos\'e V, \S2 and \S7, equation~(7.1.1).

\bibitem{EGA4}
A.~Grothendieck and J.~Dieudonn\'e,
\emph{\'El\'ements de g\'eom\'etrie alg\'ebrique IV$_3$},
Publ. Math. IH\'ES \textbf{28} (1966), (8.8.2).

\bibitem{WeilII}
P.~Deligne, \emph{La conjecture de Weil. II},
Publ. Math. IH\'ES \textbf{52} (1980), 137--252,
Th\'eor\`eme~3.3.1, as used in ZETA\_TWO ZT04--ZT06.


\bibitem{Stark} H.~M.~Stark, \emph{Some effective cases of the
Brauer--Siegel theorem}, Invent. Math. \textbf{23} (1974), 135--152.
Lemma~2 and (7), pp.~137--139; Lemma~3, pp.~139--140; Lemma~4 and
(11)--(13), pp.~140--141; Lemmas~6--7, p.~143; Theorem~3,
pp.~144--145; Lemma~8, p.~145; Lemma~9, pp.~145--146;
Theorem~$1'$ and (24)--(25), p.~148.
\url{https://gdz.sub.uni-goettingen.de/download/pdf/PPN356556735_0023/PPN356556735_0023.pdf}.
The article is printed pp.~135--152, volume PDF scans~142--159.

\bibitem{GarciaLee} S.~R.~Garcia and E.~S.~Lee,
\emph{Unconditional explicit Mertens' theorems for number fields and
Dedekind zeta residue bounds}, arXiv:2007.10313v4 (12 March 2021).
Analytic class number formula (1), p.~2; residue upper bound (2),
p.~3; \S3, pp.~16--17, explicit $c_4$ and its zero-free branch.
\url{https://arxiv.org/pdf/2007.10313}.

\bibitem{MilneANT} J.~S.~Milne, \emph{Algebraic Number Theory},
version 3.08 (2020), Theorem~4.3, printed p.~70 / PDF p.~72.
\url{https://www.jmilne.org/math/CourseNotes/ANT.pdf}.


\bibitem{NC-Louboutin2001}
S.~Louboutin, \emph{Explicit upper bounds for residues of Dedekind zeta
functions and values of $L$-functions at $s=1$, and explicit lower bounds
for relative class numbers of CM-fields}, Canad. J. Math. \textbf{53}
(2001), 1194--1222. Theorem~1(1), p.~1195; equation~(21), p.~1200.
\url{https://doi.org/10.4153/CJM-2001-045-5}.

\bibitem{NC-Wei2021}
X.~Wei, \emph{Smallness of Faltings heights of CM abelian varieties},
arXiv:2111.00617v1, 2021. Normalization and Theorem~2.1, p.~8;
Theorem~5.2, p.~19; Proposition~5.3 and proof, pp.~19--20;
Proposition~5.4 proof, p.~21.
\url{https://arxiv.org/abs/2111.00617}.

\bibitem{NC-YuanZhang2018}
X.~Yuan and S.-W.~Zhang, \emph{On the averaged Colmez conjecture},
Ann. of Math. \textbf{187} (2018), 533--638.
Normalization, p.~534; Theorem~1.1, p.~535.
\url{https://doi.org/10.4007/annals.2018.187.2.4}.

\bibitem{NC-GRPeriods2014}
\'E.~Gaudron and G.~R\'emond, \emph{Th\'eor\`eme des p\'eriodes et
degr\'es minimaux d'isog\'enies}, Comment. Math. Helv. \textbf{89}
(2014), 343--403. Theorem~1.1, p.~345; \S2.3, pp.~351--352.
\url{https://doi.org/10.4171/CMH/322}.

\bibitem{NC-GRPol2014}
\'E.~Gaudron and G.~R\'emond, \emph{Polarisations et isog\'enies},
Duke Math. J. \textbf{163} (2014), 2057--2108.
Constants, p.~2058; Theorem~1.4, p.~2060.
\url{https://doi.org/10.1215/00127094-2782528}.
Author PDF: \url{https://lmbp.uca.fr/~gaudron/art13.pdf}.

\bibitem{NC-Tsimerman2018}
J.~Tsimerman, \emph{The Andr\'e--Oort conjecture for $\mathcal A_g$},
Ann. of Math. \textbf{187} (2018), 379--390.
\S2.2 and Proposition~2.2, p.~382; Lemma~4.1, p.~384;
Theorems~4.2 and~5.2, p.~385.
\url{https://doi.org/10.4007/annals.2018.187.2.2}.

\bibitem{NC-Remond2017}
G.~R\'emond, \emph{Vari\'et\'es ab\'eliennes et ordres maximaux},
Rev. Mat. Iberoam. \textbf{33} (2017), 1173--1195.
Theorems~1.1 and~1.2; explicit same-base-field convention in the
opening of \S1.  Author preprint pp.~1--2; proof of Theorem~1.1
in Proposition~2.3, p.~6.
\url{https://www-fourier.univ-grenoble-alpes.fr/~remond/5783.pdf}.

\bibitem{NC-Binyamini2020}
G.~Binyamini, with an appendix by E.~Kowalski,
\emph{Some effective estimates for Andr\'e--Oort in $Y(1)^n$},
J. Reine Angew. Math. \textbf{767} (2020), 17--35,
doi:10.1515/crelle-2019-0028.
The locators above use arXiv:1809.05302v2 (27 May 2021),
\url{https://arxiv.org/pdf/1809.05302}.


% Bibliography entries for task 2; insert in the master thebibliography.
\bibitem{NXLBorger} J. Borger, \emph{$\Lambda$-rings and the field with one element},
arXiv:0906.3146. Supplied source
\path{LIT_Borger_Lambda_rings_and_the_field_with_one_element/lrfoe13.tex};
primary PDF edition dated 26 November 2024, pp.~2 and~6 (definition and flat
Frobenius-lift criterion), \S6 opening and Theorem~6.1, pp.~19--20
(abelian finite-coefficient Galois action), Corollary~6.2 and equation~(6.2.1),
p.~20 (potential integral-Tate splitting). Supplied source labels
\texttt{thm:lambda-varieties-are-abelian} and
\texttt{cor:lambda-varieties-are-Artin-Tate}.
\url{https://arxiv.org/pdf/0906.3146}.

\bibitem{NXLScholze} P. Scholze, \emph{Perfectoid spaces},
Publ. Math. Inst. Hautes \'Etudes Sci. \textbf{116} (2012), 245--313.
The primary author PDF used here is the 19 November 2011 version:
Definition~5.1(ii), Theorem~5.2, and Lemma~5.3, pp.~21--23;
Proposition~5.13(i) and Remark~5.14, p.~25;
Lemma~6.2 and Theorem~6.3, p.~30;
Theorem~7.12, p.~43;
Definition~7.14, pp.~43--44, and Theorem~7.17, Corollaries~7.18--7.19,
p.~44;
Theorem~9.7 (citing Huber, Theorem~3.8.1), p.~49.
\url{https://people.mpim-bonn.mpg.de/scholze/PerfectoidSpaces.pdf}.

\bibitem{NXLHuber} R. Huber, \emph{A generalization of formal schemes and rigid
analytic varieties}, Math. Z. \textbf{217} (1994), 513--551,
Remark~4.6(iv)(d), printed p.~544. The hypotheses are locally noetherian
schemes and a proper morphism; the characteristic-$p$ application in NX-L02
is descended to $\mathbb F_p[[u]]$ before base change.
\url{https://virtualmath1.stanford.edu/~conrad/Perfseminar/refs/Huberformalrigid.pdf}.

\bibitem{NXLBhatt} B. Bhatt, \emph{Lecture II: $\delta$-rings},
Columbia prismatic-cohomology lectures, Proposition~3.2, pp.~5--6
(unique strict-$p$ lift of a perfect $\mathbb F_p$-algebra).
The primary PDF text was retrieved; the direct local download returned 403.
The proof can use Scholze, Remark~5.14, alone.
\url{https://websites.umich.edu/~bhattb/teaching/prismatic-columbia/lecture2-delta-rings.pdf}.

\bibitem{NXLSilverman} J. H. Silverman, \emph{The Arithmetic of Elliptic Curves},
second edition, Graduate Texts in Mathematics 106, Springer, 2009,
Theorem~V.2.4, printed p.~143 (elliptic-curve zeta numerator).
The explicit trace $a_5=-2$ in NX-L01 is recomputed directly.
\url{https://www.math.ens.psl.eu/~obenoist/refs/Silverman.pdf}.



\bibitem{NGFaltings}
G. Faltings, \emph{Endlichkeitss\"atze f\"ur abelsche Variet\"aten
\"uber Zahlk\"orpern}, Invent. Math. \textbf{73} (1983), 349--366,
\S5, S\"atze 3--4, printed p.~360 (PDF p.~12).
\url{https://math.uchicago.edu/~drinfeld/Deligne's_conjecture_Manin_conf/Faltings_argument/Faltings.pdf}.

\bibitem{NGDeligneMilne}
P. Deligne and J. S. Milne, \emph{Tannakian Categories}, in
\emph{Hodge Cycles, Motives, and Shimura Varieties}, LNM 900,
Springer (1982), 101--228.  Corrected transcription of 15 August 2012,
with the original statement numbering: Propositions 6.5, 6.13--6.14,
6.21, 6.23; Theorem 6.25; Proposition 6.26 and Corollary 6.27.
\url{https://www.jmilne.org/math/xnotes/tc2012.pdf}.

\bibitem{NGDeligne82}
P. Deligne (notes by J. S. Milne), \emph{Hodge cycles on abelian
varieties}, in \emph{Hodge Cycles, Motives, and Shimura Varieties},
LNM 900, Springer (1982), 9--100, Proposition 2.9 and Theorem 2.11.
The updated transcription dated 1 October 2018 places these on
pp.~18--19.
\url{https://jmilne.org/math/Documents/Deligne82.pdf}.

\bibitem{NGCommelin}
J. Commelin, \emph{The Mumford--Tate conjecture for products of
abelian varieties}, Algebraic Geometry \textbf{6}(6) (2019), 650--677,
Lemma 3.8 (p.~655); Theorem 4.13 and Remark 4.14 (pp.~657--658);
\S4.16 and Lemma 4.17 (p.~658).
\url{https://content.algebraicgeometry.nl/2019-6/2019-6-028.pdf}.

\bibitem{NGMoonen}
B. Moonen, \emph{Families of Motives and the Mumford--Tate Conjecture},
Example 2.1.13 (p.~10 in the author's 52-page manuscript),
for the Mumford--Tate group of a non-CM elliptic curve.
\url{https://www.math.ru.nl/~bmoonen/Papers/FMMTC.pdf}.

\bibitem{NGNewtonThorne}
J. Newton and J. A. Thorne, \emph{Symmetric power functoriality for
holomorphic modular forms, II}, Publ. Math. IH\'ES \textbf{134}
(2021), 117--152, Theorem A (p.~117), Corollary B (p.~118).
\url{https://www.numdam.org/item/10.1007/s10240-021-00126-4.pdf}.

\bibitem{NGZarhin}
Y. G. Zarhin, \emph{The endomorphism rings of jacobians of cyclic
covers of the projective line}, arXiv:math/0103203v2 (8 September
2002), Theorem 1.1 and Examples 1.3(1), PDF p.~2; genus formula in
\S1, PDF p.~1, and differential description in Remark 3.1,
PDF pp.~7--8.
\url{https://arxiv.org/pdf/math/0103203}.


\end{thebibliography}
\end{document}
